% Template for PLoS
% Version 3.4 January 2017
%
% % % % % % % % % % % % % % % % % % % % % %
%
% -- IMPORTANT NOTE
%
% This template contains comments intended 
% to minimize problems and delays during our production 
% process. Please follow the template instructions
% whenever possible.
%
% % % % % % % % % % % % % % % % % % % % % % % 
%
% Once your paper is accepted for publication, 
% PLEASE REMOVE ALL TRACKED CHANGES in this file 
% and leave only the final text of your manuscript. 
% PLOS recommends the use of latexdiff to track changes during review, as this will help to maintain a clean tex file.
% Visit https://www.ctan.org/pkg/latexdiff?lang=en for info or contact us at latex@plos.org.
%
%
% There are no restrictions on package use within the LaTeX files except that 
% no packages listed in the template may be deleted.
%
% Please do not include colors or graphics in the text.
%
% The manuscript LaTeX source should be contained within a single file (do not use \input, \externaldocument, or similar commands).
%
% % % % % % % % % % % % % % % % % % % % % % %
%
% -- FIGURES AND TABLES
%
% Please include tables/figure captions directly after the paragraph where they are first cited in the text.
%
% DO NOT INCLUDE GRAPHICS IN YOUR MANUSCRIPT
% - Figures should be uploaded separately from your manuscript file. 
% - Figures generated using LaTeX should be extracted and removed from the PDF before submission. 
% - Figures containing multiple panels/subfigures must be combined into one image file before submission.
% For figure citations, please use "Fig" instead of "Figure".
% See http://journals.plos.org/plosone/s/figures for PLOS figure guidelines.
%
% Tables should be cell-based and may not contain:
% - spacing/line breaks within cells to alter layout or alignment
% - do not nest tabular environments (no tabular environments within tabular environments)
% - no graphics or colored text (cell background color/shading OK)
% See http://journals.plos.org/plosone/s/tables for table guidelines.
%
% For tables that exceed the width of the text column, use the adjustwidth environment as illustrated in the example table in text below.
%
% % % % % % % % % % % % % % % % % % % % % % % %
%
% -- EQUATIONS, MATH SYMBOLS, SUBSCRIPTS, AND SUPERSCRIPTS
%
% IMPORTANT
% Below are a few tips to help format your equations and other special characters according to our specifications. For more tips to help reduce the possibility of formatting errors during conversion, please see our LaTeX guidelines at http://journals.plos.org/plosone/s/latex
%
% For inline equations, please be sure to include all portions of an equation in the math environment.  For example, x$^2$ is incorrect; this should be formatted as $x^2$ (or $\mathrm{x}^2$ if the romanized font is desired).
%
% Do not include text that is not math in the math environment. For example, CO2 should be written as CO\textsubscript{2} instead of CO$_2$.
%
% Please add line breaks to long display equations when possible in order to fit size of the column. 
%
% For inline equations, please do not include punctuation (commas, etc) within the math environment unless this is part of the equation.
%
% When adding superscript or subscripts outside of brackets/braces, please group using {}.  For example, change "[U(D,E,\gamma)]^2" to "{[U(D,E,\gamma)]}^2". 
%
% Do not use \cal for caligraphic font.  Instead, use \mathcal{}
%
% % % % % % % % % % % % % % % % % % % % % % % % 
%
% Please contact latex@plos.org with any questions.
%
% % % % % % % % % % % % % % % % % % % % % % % %

\documentclass[10pt,letterpaper]{article}
\usepackage[top=0.85in,left=2.75in,footskip=0.75in]{geometry}

% amsmath and amssymb packages, useful for mathematical formulas and symbols
\usepackage{amsmath,amssymb}

% Use adjustwidth environment to exceed column width (see example table in text)
\usepackage{changepage}

% Use Unicode characters when possible
\usepackage[utf8x]{inputenc}

% textcomp package and marvosym package for additional characters
\usepackage{textcomp,marvosym}

% cite package, to clean up citations in the main text. Do not remove.
\usepackage{cite}

% Use nameref to cite supporting information files (see Supporting Information section for more info)
\usepackage{nameref,hyperref}

% line numbers
\usepackage[right]{lineno}

% ligatures disabled
\usepackage{microtype}
\DisableLigatures[f]{encoding = *, family = * }

% color can be used to apply background shading to table cells only
\usepackage[table]{xcolor}

% array package and thick rules for tables
\usepackage{array}

% create "+" rule type for thick vertical lines
\newcolumntype{+}{!{\vrule width 2pt}}

% create \thickcline for thick horizontal lines of variable length
\newlength\savedwidth
\newcommand\thickcline[1]{%
  \noalign{\global\savedwidth\arrayrulewidth\global\arrayrulewidth 2pt}%
  \cline{#1}%
  \noalign{\vskip\arrayrulewidth}%
  \noalign{\global\arrayrulewidth\savedwidth}%
}

% \thickhline command for thick horizontal lines that span the table
\newcommand\thickhline{\noalign{\global\savedwidth\arrayrulewidth\global\arrayrulewidth 2pt}%
\hline
\noalign{\global\arrayrulewidth\savedwidth}}

% Remove comment for double spacing
\usepackage{setspace} 
\doublespacing

\renewcommand*\familydefault{\sfdefault} 

% Text layout
\raggedright
\setlength{\parindent}{0.5cm}
\textwidth 5.25in 
\textheight 8.75in

% Bold the 'Figure #' in the caption and separate it from the title/caption with a period
% Captions will be left justified
\usepackage[aboveskip=1pt,labelfont=bf,labelsep=period,justification=raggedright,singlelinecheck=off]{caption}
\renewcommand{\figurename}{Fig}

% Use the PLoS provided BiBTeX style
\bibliographystyle{scripts/plos2015}

% Remove brackets from numbering in List of References
\makeatletter
\renewcommand{\@biblabel}[1]{\quad#1.}
\makeatother

% Leave date blank
\date{}

% Header and Footer with logo
\usepackage{lastpage,fancyhdr,graphicx}
\usepackage{epstopdf}
\pagestyle{myheadings}
\pagestyle{fancy}
\fancyhf{}
\setlength{\headheight}{27.023pt}
\lhead{\includegraphics[width=2.0in]{scripts/PLOS-submission.eps}}
\rfoot{\thepage/\pageref{LastPage}}
\renewcommand{\footrule}{\hrule height 2pt \vspace{2mm}}
\fancyheadoffset[L]{2.25in}
\fancyfootoffset[L]{2.25in}
\lfoot{\sf PLOS}

%% Include all macros below

\newcommand{\lorem}{{\bf LOREM}}
\newcommand{\ipsum}{{\bf IPSUM}}

%% END MACROS SECTION


\begin{document}
\vspace*{0.2in}

% Title must be 250 characters or less.
\begin{flushleft}
{\Large
\textbf\newline{
  Systematic identification of factors mediating accelerated mRNA 
  degradation in response to changes in environmental nitrogen.
} % Please use "sentence case" for title and headings (capitalize only the first word in a title (or heading), the first word in a subtitle (or subheading), and any proper nouns).
}
\newline
% Insert author names, affiliations and corresponding author email (do not include titles, positions, or degrees).
\\
Darach Miller\textsuperscript{1},
Nathan Brandt\textsuperscript{1},
David Gresham\textsuperscript{1}
\\
\bigskip
\textbf{1} Center for Genomics and Systems Biology, Department of
Biology, New York University, New York, 10003, USA.
\\
\bigskip

% Insert additional author notes using the symbols described below. Insert symbol callouts after author names as necessary.
% 
% Remove or comment out the author notes below if they aren't used.
%
% Primary Equal Contribution Note
%\Yinyang These authors contributed equally to this work.

% Additional Equal Contribution Note
% Also use this double-dagger symbol for special authorship notes, such as senior authorship.
%\ddag These authors also contributed equally to this work.

% Current address notes
%\textcurrency Current Address: Dept/Program/Center, Institution Name, City, State, Country % change symbol to "\textcurrency a" if more than one current address note
% \textcurrency b Insert second current address 
% \textcurrency c Insert third current address

% Deceased author note
%\dag Deceased

% Group/Consortium Author Note
%\textpilcrow Membership list can be found in the Acknowledgments section.

% Use the asterisk to denote corresponding authorship and provide email address in note below.
* dgresham@nyu.edu

\end{flushleft}
% Please keep the abstract below 300 words
\section*{Abstract}
Cellular responses to changing environments frequently
involve rapid reprogramming of the transcriptome.
Regulated changes in mRNA degradation rates can
accelerate reprogramming by clearing or stabilizing extant transcripts. 
%Budding yeast respond to an improvement in
%nitrogen-availability by triggering a transcriptional reprogramming
%that functions to upregulate ribosome biogenesis and repress
%alternative nitrogen-source catabolism. 
Here, we measured mRNA stability using 4-thiouracil labeling
in the budding yeast \textit{Saccharomyces cerevisiae}
during a nitrogen upshift and found that 78 mRNAs are subject
to destabilization. These transcripts include Nitrogen
Catabolite Repression (NCR) and carbon metabolism mRNAs,
suggesting that mRNA destabilization is a mechanism 
for targeted reprogramming of the transcriptome.
To explore the molecular basis of
destabilization we implemented a SortSeq approach to
screen the pooled deletion collection library
for \textit{trans} factors that mediate rapid \textit{GAP1}
mRNA repression.
We combined low-input multiplexed \underline{B}arcode sequencing 
with branched-DNA single-molecule mRNA \underline{F}ISH and 
\underline{F}luorescence-activated cell sorting (BFF)
to identify the Lsm1-7p/Pat1p complex and general mRNA
decay machinery as important for \textit{GAP1} mRNA clearance.
We also find that the decapping modulators 
\textit{EDC3} and \textit{SCD6}, translation
factor eIF4G2, and the 5' UTR of \textit{GAP1}
are factors that mediate rapid repression of \textit{GAP1} mRNA, 
suggesting that translational control may impact the 
post-transcriptional fate of mRNAs in response to 
environmental changes.

% Please keep the Author Summary between 150 and 200 words
% Use first person. PLOS ONE authors please skip this step. 
% Author Summary not valid for PLOS ONE submissions.   
\section*{Author summary}
Budding yeast rapidly increase their rate of growth upon re-feeding of
glutamine to a nitrogen-limited culture. Within the first several
minutes, the transcriptome begins to reprogram. 
We investigated changes in mRNA stability
using a 4-thiouracil label-chase experiment 
to track the extant transcriptome during the nitrogen upshift.
We found that a functionally enriched set of mRNA are destabilized,
including some members of the Nitrogen Catabolite Repression regulon.
To find genetic factors that may be involved, we used \textit{GAP1}
mRNA as a marker of this post-transcriptional regulation and
developed a method for estimating mRNA abundance using a combination
of mRNA FISH, FACS, and Barcode sequencing. We identified
the Lsm1-7p/Pat1p complex as having large differences in \textit{GAP1}
mRNA regulation, and further work identified that
the decapping modulators \textit{EDC3} and \textit{SCD6}, the
translation initiation factor \textit{eIF4G2},
and the 5' UTR of \textit{GAP1} have defects in
rapidly clearing \textit{GAP1} mRNA during this transition. 
Our study suggests that changes in the translational status of
specific mRNAs may impact their stability in response to environmental
signals.
%
%However,
%the phenotype of deletions of \textit{SCD6} and \textit{eIF4G2}
%suggest this may result from a relative stabilization before the shift.

\linenumbers

% Use "Eq" instead of "Equation" for equation citations.
\section*{Introduction}

Regulated changes in mRNA abundance are a primary cellular response
to external stimuli.
Both the rate of synthesis and the rate of degradation determine the
steady-state abundance of a particular mRNA and the kinetics
with which abundance changes occur
\cite{Hargrove1989,Perez-Ortin2013}. 
Changes in mRNA degradation rates fulfill an important 
mechanistic role in diverse systems, including 
development \cite{Alonso2012,West2017} and disease
\cite{Aghib1990}.
In budding yeast, the rate of
mRNA degradation is affected by environmental stresses
\cite{Canadell2015}, cellular growth rate
\cite{Garcia-Martinez2016}, as well as by improvements in 
nutrient conditions \cite{Scheffler1998}.

Environmental shifts trigger rapid reprogramming of the budding yeast
transcriptome in response to stresses and nutritional
changes \cite{Gasch2000,Conway2012}. Changes in mRNA degradation rates
have been shown to play a role in responses to heat-shock, osmotic
stress, pH increases, and oxidative stress
\cite{Castells2011,Romero2009,Canadell2015,Molina2008}. 
In response to these
diverse stresses destabilization of mRNAs encoding 
ribosomal-biogenesis gene products, and  
stress-induced mRNA occurs \cite{Canadell2015}. 
Simultaneous increases in both synthesis and
degradation rates of some  mRNAs may serve to speed the return to a
steady-state following a transient pulse of regulation
\cite{Shalem2008}. Addition of glucose to carbon-limited cells 
results in both stabilization of 
ribosomal protein mRNAs \cite{Yin2003} and destabilization
of gluconeogenic transcripts \cite{De_la_Cruz2002,Mercado1994}.
Destabilization of transcripts can
have a delayed effect on reducing protein levels compared to
up-regulated genes \cite{Lee2011}. This suggests that accelerated
mRNA degradation may serve additional purposes. For example, clearance
of specific mRNAs could increase nucleotide pools
\cite{Kresnowati2006} or facilitate reallocation of
translational capacity 
\cite{Kief1981,Giordano2016,Shachrai2010}. 
%Identifying the genetic
%factors responsible for the accelerated mRNA degradation would allow
%us to test if regulated destabilization of specific transcripts is
%adaptive.

Yeast cells metabolize a wide variety of nitrogen sources, but
preferentially assimilate and metabolize specific nitrogen compounds.
Transcriptional regulation, known as
“nitrogen catabolite repression” (NCR) \cite{Magasanik2002},
controls the expression of mRNAs
encoding transporters, metabolic enzymes, and regulatory
factors required for utilization of alternative nitrogen sources. 
NCR-regulated transcripts are expressed in the
absence of a readily metabolized (preferred) nitrogen sources or in
the presence of growth-limiting concentrations (in the low μM range)
of any nitrogen source \cite{Godard2007,Airoldi2016}. Regulation
of NCR targets is mediated by two activating GATA
transcription factors, Gln3p and Gat1p, and two repressing
GATA factors, Dal80p and Gzf3p. \textit{GAT1}, \textit{GZF3}, and
\textit{DAL80} promoters
contain GATAA motifs, and thus transcriptional regulation of NCR
targets entails self-regulatory and cross-regulatory loops. When
supplied with a preferred nitrogen source such as glutamine, the
NCR-activating transcription factors Gat1p and Gln3p are excluded from
the nucleus by TORC1-dependent and -independent mechanisms
\cite{Beck1999,Tate2013,Tate2017} and NCR transcripts are strongly
repressed. The activity of some NCR gene
products is also controlled by post-translational mechanisms
\cite{Cooper1983} such as the General Amino-acid Permease
(Gap1p) which is rapidly inactivated upon a nitrogen 
upshift via ubiquitination
\cite{Stanbrough1995,Risinger2006,Merhi2012}. Recently, we have
identified an additional level of regulation of NCR transcripts: cells
growing in NCR de-repressing conditions accelerate the degradation
of \textit{GAP1} %and \textit{DIP5} mRNAs
mRNA upon addition of glutamine
\cite{Airoldi2016}. Thus, mRNA degradation rate regulation may be an
additional mechanism for clearing NCR-regulated transcripts upon 
improvements in environmental nitrogen availability.

Multiple pathways mediate the degradation of mRNAs. The main pathway
of mRNA degradation occurs by deadenylation and decapping
prior to 5’ to 3’ exonucleolytic degradation by Xrn1p; however,
transcripts are also degraded 3’ to 5’ via the exosome, or via
activation of co-translational quality control mechanisms
\cite{Parker2012}. Deadenylation of mRNAs by the Ccr4-Not complex
allows the mRNA to be bound at the 3' end by the 
Lsm1-7p/Pat1p complex, a heptameric
ring comprising the SM-like proteins Lsm2-7p and the
cytoplasmic-specific Lsm1p \cite{Tharun2000,Sharif2013}, which then
recruits factors for decapping by Dcp2p. 
Recruitment of the decapping enzyme \cite{Coller2004} is the 
rate-limiting step for canonical 5'-3' degradation.
Therefore Lsm1-7p, Pat1p,
and associated factors play a key role in determining the kinetics
of mRNA degradation \cite{Nissan2010}. 

Regulation of mRNA degradation pathways can alter the stability of
specific mRNAs. For example, the RNA-binding protein (RBP) Puf3p
recognizes a \textit{cis}-element in 3’ UTRs \cite{Olivas2000} 
and affects mRNA degradation rates depending on
Puf3p phosphorylation status \cite{lee2015}. 
%Transcript properties
%also associated with translation dynamics affect mRNA degradation, at
%the level of elongation \cite{Sweet2012,Presnyak2015,Neymotin2016} or
%competition between the decapping enzymes and translation initiation
%\cite{Schwartz2000}. 
In addition to \textit{cis}-elements within the transcript, 
promoters have
been shown to mark certain RNA-protein (RNP) complexes to specify
their post-transcriptional regulation
\cite{Mercado1994,Haimovich2013,Trcek2011,braun2016snf1}. These
mechanisms may be controlled by a variety of different signalling
pathways including Snf1 \cite{Young2012,Braun2014}, PKA
\cite{Ramachandran2011}, Phk1/2 \cite{Luo2011}, and TORC1
\cite{Talarek2010}. Thus, regulated changes in  mRNA degradation
rates entails numerous mechanisms that collectively tune stability of
mRNAs in response to the activity of signalling pathways. 

Here, we studied the global regulation of mRNA degradation rates upon
an increase in nitrogen availability using 4-thiouracil (4tU) 
label-chase and RNAseq.
We found that a set of 78 mRNAs are subject to accelerated mRNA
degradation, including many NCR transcripts as well as mRNAs
encoding components of
carbon metabolism. To identify the mechanism underlying accelerated
mRNA degradation we designed a high-throughput genetic screen using 
\underline{B}arcode-sequencing of a pooled library which was
fractionated using \underline{F}luorescence-activated cell 
sorting on the basis of single molecule mRNA \underline{F}ISH signal (BFF). 
We screened the barcoded
yeast deletion collection to test the effect of each gene deletion
on the abundance of \textit{GAP1} mRNA in NCR de-repressing 
conditions and its clearance following the 
addition of glutamine. We
find that the Lsm1-7p/Pat1p complex and decapping modifiers affect
both \textit{GAP1} mRNA steady-state expression and its 
accelerated degradation.
This work expands our
understanding of mRNA stability regulation in remodeling the
transcriptome during a relief from growth-limitation and demonstrates
a generalizable approach to the study of genetic determinants of mRNA
dynamics.

\section*{Results}

\subsection*{Transcriptional reprogramming precedes physiological remodeling}

Cellular responses to environmental signals entail coordinated changes
in both gene expression and cellular physiology.  Previously, we
studied the steady-state and dynamic responses of 
\textit{Saccharomyces cerevisiae} 
(budding yeast) to environmental nitrogen
\cite{Airoldi2016}, and found that the transcriptome is rapidly
reprogrammed following a single pulsed addition of glutamine to
nitrogen-limited cells in either a chemostat or
batch culture. To study physiological changes in response to a
nitrogen upshift, we measured growth rates of a population of 
cells. A prototrophic haploid lab strain 
(FY4, isogenic to S288c) grows with a
4.5 hour doubling time in batch culture in minimal media 
containing proline as a sole
nitrogen source (Fig~\ref{figure1}A). Upon addition of 400μM glutamine
the cells undergo a 2-hour lag period during which no change in
population growth rate is detected, but the average cell size
continuously increases ($\sim$21\% increase in mean volume 
Fig~\ref{figure1}B). Following the 2-hour lag, the population adopts a 2.1 
hour doubling time.
%This lag in population growth rate upon an upshift has been 
%described before \cite{Carter1978}.
By contrast, global gene expression changes are detected
within three minutes of the upshift \cite{Airoldi2016}. 
Thus, transcriptome remodeling precedes
physiological remodeling in response to a nitrogen upshift.

To evaluate concordance in transcriptome remodeling between chemostat
and batch nitrogen upshifts, and the extent to which they reflect
changes in gene expression observed during systematic steady-state 
changes in growth rates using chemostats, we
performed principal component analysis of global gene expression
(Fig~\ref{figure1}C). The first two principal components, which
account for almost half of the total variation, clearly separate
steady-state and nitrogen upshift cultures.  Systematic changes in
growth rate primarily results in
separation of gene expression states along the second principal
component, whereas upshift experiments vary along the first 
principal component.  This suggests that
although a nitrogen upshift results in a gene expression state 
reflecting increased cell growth rates \cite{Airoldi2016}, the
transcriptome is remodeled through a distinct state. 
In upshift experiments in
chemostats, the gene expression trajectory ultimately returns to 
the initial steady-state condition as excess nitrogen is 
depleted by consumption and dilution 
(\nameref{longTermPCA}, \nameref{timeplot}). 

To investigate the functional basis of gene expression programs
in the upshift and steady-state conditions, we computed the
correlation of each transcript with the loadings on these first two
principal components and performed gene-set enrichment analysis
(\nameref{microarrayPCAgsea}). 
Component 1 is positively correlated with functions including 
mRNA processing, transcription from RNA polymerases (I,II,and III),
and chromatin organization, and negatively correlated with
cytoskeleton organization,
vesicle organization, membrane fusion, and cellular respiration.
Both steady-state and upshift gene expression
trajectories increase with principal component 2, but they diverge
along principal component 1. Components 1 and 2 are 
strongly enriched for terms including ribosome biogenesis, 
nucleolus, and
tRNA processing, and negatively correlated with
vacuole, cell cortex, and carbohydrate metabolism terms. 
Together, this analysis suggests that both upshift and
increased steady-state growth rates share upregulation of
ribosome-associated components, but the reprogramming
preceeding the upshift in growth reflects an immediate focus on 
quantitative changes in
gene expression machinery instead of structural cellular components.
Importantly,
dynamic reprogramming is similar in both the chemostat and batch
upshift (Fig~\ref{figure1}C). As batch cultures are a technically
simpler experimental system, we performed all subsequent experiments
using batch culture nitrogen upshifts. 

\subsection*{Global analysis of mRNA stability changes during the
nitrogen upshift}

Previously, we found that \textit{GAP1} and \textit{DIP5} mRNAs 
are destabilized in
response to a nitrogen upshift \cite{Airoldi2016}. We sought to
determine if mRNA destabilization is specific to NCR transporter
mRNAs by measuring global mRNA stability across the upshift
using 4-thiouracil (4tU) labeling and RNA-seq 
\cite{Neymotin2014,Munchel2011}.
As a 4tU labeling experiment requires uracil transport, 
which may be altered upon stresses or a 
change in nitrogen-availability \cite{volland1994,Hein1995}, 
we designed experiments such
that %following complete 4tU labeling and metabolism to nucleotides
the chase was initiated prior to addition of glutamine or water (mock).
We normalized data using \textit{in vitro} synthesized thiolated 
spike-ins by fitting a log-linear model to spike-in counts
across time (\nameref{writeup2}), which reduced noise and increased
our power to detect stability changes (\nameref{dme211raw},
\nameref{dme211filterDirect},
\nameref{dme211filterModel}).
Data and models for each transcript can be visualized in browser
using a Shiny appplication (
\url{http://shiny.bio.nyu.edu/users/dhm267/}). 

We modeled the
log-transformed normalized signal for each mRNA using linear
regression (\nameref{dme211resultsModel}).
Of 4,859 mRNAs measured we identified 94 that increased in 
degradation rate and 38 that decreased (FDR $<$ 0.01, using \cite{Storey2003}).  %inlineciteStorey2003
We generated a model of nucleotide
labeling kinetics to assess the effect of an incomplete label 
chase on our experimental results (\nameref{writeup2}),
 and found that complete transcriptional inhibition alone could 
theoretically result in a 13.3\% increase in the apparent 
degradation rate. In order to eliminate the possibility that
rapid synthesis changes could affect our estimates,
we considered transcript destabilization to be at least a
doubling (100\% increase) in the apparent degradation rates between 
pre-upshift and post-upshift.
This conservative cutoff 
left 78 mRNA that are significantly destabilized 
upon a nitrogen upshift. 

The vast majority of transcripts (4,781 of 4,859) do not show
evidence for stability changes upon addition of glutamine
(e.g. \textit{HTA1}, Fig~\ref{figure2}A). 
The median pre-upshift half-life is 6.92 minutes and the median
half-life following the upshift is 6.32 minutes (Table~\ref{table1})
suggesting that there is not a global change in mRNA stability.
Global stability estimates are
considerably lower than previous estimates in rich medium
\cite{Munchel2011,Neymotin2014,Miller2011}, which may reflect the
different nutrient conditions used in our study. 
The 78 transcripts significantly destabilized upon the 
glutamine-upshift include
mRNAs encoding NCR transporters \textit{GAP1}, \textit{DAL5}, and
\textit{MEP2} (blue label, Fig~\ref{figure2}A), the pyruvate metabolism enzymes
\textit{PYK2} and \textit{PYC1} (orange label), and trehalose synthase
subunits \textit{TPS1} and
\textit{TPS2} (yellow label).
Destabilized mRNA tend to be more stable before the upshift (Fig~\ref{figure2}B),
(median half-life of 9.46 minutes) and exhibit 
a median 3.06-fold increase in degradation rates (median half-life of
3.02 minutes following the upshift). 

\begin{table}[!ht]
\begin{adjustwidth}{-2.25in}{0in} % Comment out/remove adjustwidth environment if table fits in text column.
\centering
\caption{\textbf{Summary of mRNA stability and changes upon the upshift.} Shown here are the median rates or changes in rates for the specified sets. Destabilized transcripts were identified using the criteria of a significant (FDR $<$ 0.01) change in estimated degradation rates and at least a doubling of the rate of clearance.}
\begin{tabular}{p{8em} | l l l l | l l}
& \multicolumn{2}{c}{Pre-shift} & \multicolumn{2}{c}{Post-shift} &
Change in & Fold-change\\
 & specific rate & half-life & specific rate & half-life & specific
rate & specific rate\\
 & (min$^{-1}$) & (min) & (min$^{-1}$) & (min) & (min$^{-1}$) & \\
\\
\raggedright All transcripts & 0.100 & 6.92 & 0.110 & 6.32 & 0.00865 &
1.08\\
\raggedright Destabilized (n=78) & 0.0732 & 9.46 & 0.229 & 3.02 &
0.158 & 3.06\\
\raggedright No difference\\detected (n=4781) & 0.101 & 6.89 & 0.108 &
6.40 & 0.00728 & 1.07\\
\end{tabular}
\label{table1} 
\end{adjustwidth}
\end{table}

We tested for
functional enrichment among the set of 78 destabilized
mRNAs and found that they are strongly enriched for NCR
transcripts (16 of 78, p $<$ $10^{-11}$). Almost half of the
destabilized transcripts are annotated as “ESR-up” genes
(\nameref{comparisonESR}), on the basis of  their rapid induction
during the environmental stress response \cite{Gasch2000}. 78
destabilized mRNAs are enriched (FDR $<$ 0.05) for additional GO terms and KEGG 
pathways (\nameref{dme211goAndKegg}) including
glycolysis/gluconeogenesis (6/78 genes), 
carbohydrate metabolic process (24/78 genes),
trehalose-phosphatase activity (3/78 genes), 
pyruvate metabolic process (6/78 genes), 
and secondary active transmembrane transport
(8/78 genes, a subset of the enriched 11 ion transmembrane transport genes).
%We also see destabilization of \textit{PYK2} and \textit{HXK1},
%both of which are isozymes expressed highly in poor nutrient conditions.
Thus destabilization upon a nitrogen upshift regulates, 
but is not restricted to, NCR-regulated mRNAs and reflects broader
metabolic changes in the cell. 





To investigate the extent to which mRNA stability changes contribute
to transcriptome reprogramming, we compared degradation rates
to abundance changes\cite{Airoldi2016} following the upshift 
(Fig~\ref{figure2}C). 
Changes in mRNA degradation rates %inlineciteAiroldi2016
and expression change rates are anti-correlated (Pearson's $r$ = -0.598,
p-value $<$ $10^{-15}$, \nameref{kkdComparison}),
consistent with stability changes impacting gene expression dynamics.
However, they are not entirely co-incident, as some destabilized
transcripts do not exhibit decreases in abundance (red points in
Fig~\ref{figure2}C, \nameref{comparisonDestabilized},
and \nameref{compareSix}).
This analysis indicates that increases in degradation rates play a 
significant role
in the rapid reprogramming of the transcriptome upon a glutamine
upshift, but that in some cases cases they are counteracted by
increases in mRNA synthesis rates \cite{Shalem2008,Canadell2015}. 

Functional coordination of mRNA stability changes suggests  a possible
role for \textit{cis}-element regulation. We analyzed UTRs and coding
sequence for enrichment of new motifs or known RNA binding protein
(RBP) motifs.
3’ UTRs of destabilized transcripts are
depleted of Puf3p binding sites, but we found no enriched sequence
motif in the 3' UTRs.
5’ UTRs are enriched for a GGGG motif, which
may be explained by convergence between mRNA stability changes and
transcriptional control by Msn2/4 on the ESR “up” genes
(\nameref{comparisonESR}, \cite{Gasch2000,Canadell2015}).  %inlineciteGasch2000,Canadell2015
5’ UTRs are also enriched for binding motifs reported for Hrp1p 
(\nameref{hrp1}),
a canonical member of the nuclear cleavage factor I complex \cite{Chen1998}.
However, this protein has been shown to shuttle to the cytoplasm
where it may play a regulatory role
\cite{Kessler1997,Kebaara2003,Guisbert2005}.
On average,
destabilized mRNAs are longer and contain more optimal codons
(\nameref{lengthAndCodons}, \cite{Khong2017}).  %inlineciteKhong2017
Together, this analysis suggests that the
mechanism of destabilization may act through cis elements in the 5’
UTR and/or biased codon usage.

\subsection*{A genome-wide screen for \textit{trans}-factors regulating \textit{GAP1} mRNA repression}

We sought to identify \textit{trans}-factors mediating accelerated mRNA
degradation in response to a nitrogen upshift. We selected \textit{GAP1} 
as representative of transcript destabilization, as it is abundant in
nitrogen-limiting conditions and is rapidly cleared upon addition of
glutamine  (3.24-fold increase in degradation rate, Fig~\ref{figure3}A,
\nameref{dme211resultsModel}). Previous approaches to high-throughput
genetics of transcriptional activity have used protein expression
reporters \cite{Neklesa2009,Sliva2016} or automation of qPCR 
\cite{Worley2015}. However, for our
purposes, we required direct measurement of \textit{GAP1} mRNA 
changes on a rapid timescale.
Therefore, we applied 
\underline{s}ingle 
\underline{m}olecule 
\underline{f}luorescent 
\textit{\underline{i}n \underline{s}itu}
\underline{h}ybridization (smFISH) to quantify 
native \textit{GAP1} transcripts in yeast cells in the pooled
prototrophic yeast deletion collection \cite{Vandersluis2014}.
Using fluorescence activated cell sorting (FACS) and Barseq
\cite{Smith2009,Robinson2013,Giaever2014},
we aimed to quantify and model the distribution of \textit{GAP1} mRNA
in each mutant \cite{Kinney2010,Peterman2016}.

%Development of our screen required that we could detect and
%sort cells using \textit{GAP1} mRNA signal. 
We found that
individually labeled probes tiled across \textit{GAP1} mRNA
\cite{Raj2008} were insufficiently bright for
\textit{GAP1} mRNA quantification using flow cytometry,
likely due to the small cell size of nitrogen-limited cells and the
low transcript numbers in yeast cells compared to mammalian cells
\cite{Klemm2014}. Therefore, we used branched DNA probes
(Quantigene), which serve to amplify the FISH signal
\cite{Hanley2013}. We developed a fixation and permeabilization
protocol (\nameref{writeup4}) that enabled detection of the
distribution of  \textit{GAP1} mRNA in steady-state nitrogen-limited conditions
and its repression following the  upshift (Fig~\ref{figure3}B). In control
experiments, we found that the signal is eliminated in a \textit{GAP1} deletion strain
or by omitting the targeting probe % that confers specificity
(Fig~\ref{figure3}B and \nameref{gap1Delete}). To validate
sorting, we sorted a sample of cells into quartiles and used
microscopy to count fluorescent foci per cell
(Fig~\ref{figure3}C).
We found that increased flow cytometry signal is associated with an
increase in the number of foci in the cells (Fig~\ref{figure3}D, $R^2$ = 0.607,
p $<$ $10^{-11}$). 

Previous SortSeq studies of
the yeast deletion collection have used outgrowth 
to generate sufficient material for 
Barseq \cite{Sliva2016}. However, formaldehyde fixation precludes
outgrowth. We found that below approximately $10^6$ templates, the
Barseq reaction produces primer dimers
that outcompete the intended PCR product (\nameref{writeup4}). 
Therefore, we re-designed the
PCR reaction \cite{Robinson2013,Smith2009} to be robust for
low sample inputs (\nameref{writeup4}). Our protocol
incorporates a 6-bp \underline{u}nique \underline{m}olecular 
\underline{i}dentifier (UMI) into the first
round of extension to identify PCR duplicates, 
and uses 3’-phosphorylated oligonucleotides and a
strand-displacing polymerase (Vent exo-) to block primer dimer formation and 
off-target amplification. 
%We developed a bioinformatics pipeline 
%using pairwise alignment
%for per-read quality-filtering and compatibility with variable barcode
%length, and using the degenerate UMI barcodes to help account for PCR
%duplicates. 
%UMIs to identify duplicates.
Because strain barcodes are of variable lengths, 
we developed a bioinformatic pipeline to extract barcodes and UMIs 
using pairwise alignment to invariant flanking sequences.
Based on \textit{in silico} benchmarks, this
approach was robust to systematic and simulated random errors 
that can confound analysis of the yeast deletion barcodes 
(\nameref{codeanddata}, \nameref{writeup4}). 

We refer to this experimental approach as BFF 
(\underline{B}arseq after \underline{F}ACS after \underline{F}ISH). 
We used BFF to estimate \textit{GAP1} mRNA abundance for every mutant in the
haploid prototrophic deletion collection \cite{Vandersluis2014} in
nitrogen-limiting conditions and 10 minutes following the upshift. 
This approach facilitates identification of mutants with
defects in mRNA regulation at both the transcriptional and
post-transcriptional level without altering \textit{GAP1} mRNA 
\textit{cis}-elements that may affect its regulation. 
Moreover, this design enables identification of factors that 
regulate both the steady-state abundance of \textit{GAP1} mRNA and 
its transcriptional repression following an upshift.
We analyzed the deletion pool in biological triplicate
(Fig~\ref{figure4}A). Following barcode sequencing we found that UMIs 
approached saturation at a slower rate than expected for random sampling,
consistent with PCR amplification bias 
(\nameref{rarefaction}), and therefore we used an error
correction model \cite{Fu2011}. After 
filtering, we calculated a
pseudo-events metric that approximates the number of each mutant sorted
into each bin. 
Principal components analysis shows that the samples are 
separated primarily by FACS bin within each
condition and biological replicates are clustered indicating that our
approach reproducibly captures the variation of  \textit{GAP1} mRNA flow
cytometry signal across the library (\nameref{pca}). 
 
\subsection*{Estimating \textit{GAP1} mRNA abundance for individual mutants}

We estimated the distribution of \textit{GAP1} mRNA for each mutant by
modeling pseudo-events in each quartile as a
log-normal distribution using likelihood maximization  
(Fig~\ref{figure4}B). 
From model fits we estimated the mean expression value for each
mutant and found that the distribution of means estimated for
3,230 strains (\nameref{dme209pooledFits}, Fig~\ref{figure4}C) 
recapitulates the overall
distribution of flow cytometry signal (Fig~\ref{figure4}A). 
%Specifically, replicate A had a consistently lower estimate of
%\textit{GAP1} FISH fluoresence in both flow cytometry and modeling.
%Replicate C had fewer mutants sorted (\nameref{writeup4}), 
%reflected in the wider distribution of estimated means.
%To estimate \textit{GAP1}
%mRNA per strain, we used all replicate measurement to perform model
%fitting and filtered models for sufficient measurements  (at least two
%of three replicates in at least three of the four bins). We generated
%expression distribution estimates for 3,230 strains, and used the mean
%of each distribution as the estimate of \textit{GAP1} mRNA abundance for each
%strain (\nameref{dme209pooledFits}). %added after commented out
To validate our approach we first examined
strains for which we expected to have a specific phenotype and
compared their mean expression level to the distribution of expression
for the entire population (Fig~\ref{figure4}D). We found that the wildtype
genotype (\textit{his3}$\Delta$ complemented by the spHis5 allele during
library construction) has an expression level that is centrally
located in the distribution both before and following the upshift. The
\textit{gap1}$\Delta$ genotype is a negative control and 
we estimate that it is at the extreme
low end of the distribution before and following the upshift. 
Dal80p is a direct transcriptional repressor
of NCR transcripts %like \textit{GAP1}, 
and we found that the \textit{dal80}$\Delta$ genotype is defective in
repression of \textit{GAP1} before and after the upshift. 
Counter-intuitively, deletion of \textit{GAT1}, a transcriptional activator
of \textit{GAP1}, appears to have higher steady-state expression of
\textit{GAP1} mRNA.
However, increased expression of \textit{GAP1} mRNA in a
\textit{gat1}$\Delta$ background has
previously been reported \cite{Scherens2006} and is thought to
result from the complex interplay of NCR transcription factors on
their own expression levels. 
Data and models for each mutant strain can be visualized in browser
using a Shiny appplication (
\url{http://shiny.bio.nyu.edu/users/dhm267/} ). 

To identify new cellular processes that regulate \textit{GAP1} mRNA abundance, we
used gene-set enrichment analysis (\nameref{dme209gsea}).
Following the upshift we found mutants that
maintain high \textit{GAP1} mRNA expression are enriched for negative
regulation of gluconeogenesis (\nameref{gluco}) and the
Lsm1-7p/Pat1p complex (Fig~\ref{figure5}A). Mutants in the TORC1 signalling
pathway were not enriched; 
however, we found that a \textit{tco89}$\Delta$ mutant has
greatly increased \textit{GAP1} mRNA expression before and after the upshift
(\nameref{tco89}), consistent with the repressive role of TORC1
on the NCR regulon.
To compare expression before and after the upshift for each mutant,
we regressed the post-upshift mean expression against the pre-upshift 
mean expression for each genotype (\nameref{prePredictPost}). 
We used the residuals for each
strain to identify mutants that clear \textit{GAP1} mRNA with kinetics slower
than expected by this model.
We found that the Lsm1-7p/Pat1p complex is again strongly 
enriched for slower than
expected \textit{GAP1} mRNA clearance (\nameref{dme209pooledFits}). 
Specifically,
the \textit{lsm1}$\Delta$, \textit{lsm6}$\Delta$, and 
\textit{pat1}$\Delta$ strains have highly elevated \textit{GAP1}
mRNA expression before the upshift and are strongly impaired in the 
repression of \textit{GAP1} mRNA after the upshift (Fig~\ref{figure5}A). 

As these factors are associated with processing-body dynamics, 
we tested if microscopically-observable processing-bodies form or
disassociate during the upshift, using microscopy of Dcp2-GFP. 
We did not observe qualitative changes
in Dcp2-GFP distribution (\nameref{pbodyScope}),
and thus the upshift does not
result in a microscopically visible changes in processing-body foci.
%as seen in other stresses. 
This is consistent with previous
investigations of amino-acid limitation stress \cite{Hoyle2007} and
suggests that the defects in \textit{GAP1} mRNA clearance 
in mutants defective in the Lsm1-7p/Pat1p complex
likely result from their roles in decapping or associated processes.



To confirm the role of the Lsm1-7p/Pat1p  complex in clearing \textit{GAP1}
mRNA during the nitrogen upshift we measured \textit{GAP1} mRNA
repression using qPCR measurements normalized to
\textit{HTA1}, which is not subject to destabilization upon the upshift
(Fig~\ref{figure2}A). We also tested mutants that were not detected using BFF,
or were only detected in the highest \textit{GAP1} bin and therefore
not suitable for modeling
(e.g. \textit{xrn1}$\Delta$ \nameref{tco89}). 
Using this assay we found that the main 5’-3’ 
exonuclease \textit{xrn1}$\Delta$ 
and mRNA deadenylase complex (\textit{ccr4}$\Delta$ and
\textit{pop2}$\Delta$) are impaired in \textit{GAP1} mRNA repression 
(Fig~\ref{figure5}B).
We also found that qPCR confirmed results from the BFF assay.
We confirmed that the accelerated degradation of \textit{GAP1} mRNA is impaired
in \textit{lsm1}$\Delta$ and \textit{lsm6}$\Delta$ 
(Fig~\ref{figure5}C). 
We tested
\textit{scd6}$\Delta$ and \textit{edc3}$\Delta$, two modifiers of the
decapping or processing-body
assembly functions associated with this complex, and found two
distinct phenotypes (Fig~\ref{figure5}D). \textit{edc3}$\Delta$ has similar expression 
as wild-type before the upshift, but clears \textit{GAP1} mRNA much more slowly.
\textit{scd6}$\Delta$ has greatly reduced \textit{GAP1} mRNA expression
before the upshift but is impaired in \textit{GAP1} mRNA clearance. 
Interestingly, \textit{tif4632}$\Delta$, a deletion of eIF4G2
known to interact with Scd6p \cite{Rajyaguru2012}, 
exhibits a similar phenotype. 

Identification of an initiation factor subunit mutant with defects in
\textit{GAP1} mRNA clearance suggests that translation control may
impact stability changes. Therefore we deleted the 5' UTR
and 3' UTR of \textit{GAP1}. 
%100bp and 152bp upstream of the start codon (approximate 5’ UTR) or the
%100bp downstream of the stop codon (approximate 3’ UTR), 
Whereas the 3’ UTR deletion does not have an effect, the 5’ UTR deletion
exhibits the phenotype of reduced \textit{GAP1} mRNA before the upshift
and a reduced rate of transcript clearance following the upshift
(Fig~\ref{figure5}E). 
We observed a similar phenotype with an independent deletion of 
152bp upstream of the \textit{GAP1} start codon (\nameref{bothutr}). 
%These measurements suggest altered mRNP composition of the
%Lsm1-7p/Pat1p complex and associated decapping factors are associated
%with defects in \textit{GAP1} mRNA expression dynamics upon a 
%nitrogen upshift. Importantly the phenotype of the \textit{scd6}$\Delta$, 
%\textit{tif4632}$\Delta$, or 5’ sequence deletions preceed the 
%addition of glutamine, suggesting that the observed destabilization
%of \textit{GAP1} may be the halt of a stabilization effect, perhaps
%due to changes in translational status of \textit{GAP1}.
These results indicate that \textit{cis}-elements responsible for the
rapid clearance of \textit{GAP1} mRNA are unlikely to be located in the
3' UTR and instead may be present in the 5' end of the
mRNA.

\section*{Discussion}

Regulated changes in mRNA stability allows cells to rapidly reprogram
gene expression, clearing extant transcripts that are no longer
required and potentially reallocating translational capacity.
%Despite progress in understanding the pathways that mediate
%mRNA degradation, the functional role of mRNA degradation and the
%factors that control regulated changes in mRNA stability remain poorly
%understood. 
Pioneering work in budding yeast has shown that mRNA
stability changes facilitate gene expression remodeling in response to
changes in nutrient availability including changes in carbon sources
\cite{Scheffler1998} and iron starvation \cite{Puig2005}. 
Here, we characterized genome-wide changes
in mRNA stability in response to changes in nitrogen availability and
identified factors that mediate the rapid repression of the
destabilized mRNA, \textit{GAP1}. Our study extends our previous work
characterizing the dynamics of transcriptome changes using chemostat
cultures \cite{Airoldi2016} and shows that accelerated mRNA
degradation targets a specific subset of the transcriptome in response
to changes in nitrogen availability. We developed a novel approach to
identify regulators of mRNA abundance using pooled mutant screens and
find that modulators of decapping activity, and core degradation
factors, are required for accelerated degradation of 
\textit{GAP1} mRNA. 
 
Measuring the stability of the transcriptome requires the ability to
separate pre-existing and newly synthesized transcripts. We modified
existing methods to measure 
post-transcriptional regulation of the yeast transcriptome in a
nitrogen upshift using 4-thiouracil labeling
\cite{Miller2011,Neymotin2014,Munchel2011}. These
modifications entailed improved normalization and quantification of
extant transcripts and explicit modeling of labelling dynamics to
account for some of the inherent limitations of metabolic labeling
approaches \cite{Perez-Ortin2013}. Continued development of
fractionation biochemistry \cite{Duffy2015}  and incorporation of
explicit per-transcript efficiency terms will improve these
methods further \cite{Chan2017}.

Our experiments show that the process of physiological and gene
expression remodeling occur on very different timescales in response
to a nitrogen upshift. Cellular physiology is remodeled over the
course of two hours to achieve a new growth rate.
By contrast, transcriptome remodeling occurs rapidly and through
states that are distinct from 
those observed during systematic
increases in steady-state growth rates. 
%Interestingly,
%we found that the yeast transcriptome is on average less stable in
%both nitrogen poor conditions (6.89 min) and following the upshift
%(6.40 min) compared to rich media conditions reported in other studies
%using metabolic labeling
%\cite{Munchel2011,Neymotin2014,Miller2011}. This relative
%reduction in mRNA stability could be an adaptation to potentially
%limiting ribonucleotides, but further work exploring differences in
%mRNA degradation rates during growth limited by different nutrients is
%required to test this concept \cite{Garcia-Martinez2016}.
%Stability changes upon the nitrogen upshift generally exhibited the
%expected relationship with rates of abundance change (
%anti-correlation, $R^2=$-0.376 ); however, we found multiple cases in
%which increased mRNA degradation rates did  not result in rapid
%decreases in mRNA abundance. This has been observed for transcripts
%up-regulated in stress conditions, and has been proposed as a
%mechanism to effect a rebalancing of the transcriptome after a
%transient phase of reprogramming \cite{Shalem2008}. Importantly,
%the changes in mRNA stability that we detect are nearly coincident
%with the environmental perturbation suggesting that a signal is sensed
%and the effect propagated to impact post-transcriptional regulation
%with rapid kinetics.
We have previously shown the rapid dynamics of NCR regulon repression 
upon a nitrogen upshift \cite{Airoldi2016}. Our results here
indicate that accelerated degradation of 
%at least 16 of the 77 probable 
many NCR transcripts \cite{Godard2007} contributes to this
repression. 
%A three-fold increase in
%the degradation rate of \textit{GAP1} mRNA provides an additional layer of
%repressive control. 
Importantly, our results show that accelerated
degradation is not limited to NCR transcripts but also targets
transcripts enriched in carbon metabolism pathways, particularly
pyruvate metabolism. Conversely, we also detect an apparent reduction in the 
degradation rate for some transcripts 
%enriched in ribosome biogenesis mRNAs (for example \textit{NSR1}) as well as 
including \textit{MAE1}. \textit{MAE1} encodes
an enzyme responsible for the conversion of malate to pyruvate, and
combined with the accelerated degradation of \textit{PYK2} mRNA 
may reflect the adoption of an adaptive shunt of carbon skeletons from glutamine 
to alanine via the TCA cycle \cite{Boles1998}. 
%Due to the limitations of
%labeling approaches \cite{Perez-Ortin2013} we cannot conclude 
%here that these transcripts are indeed stabilized, 
%however we can conclude that they are strongly upregulated. 
A recent study %inlineciteTesniere2017
described destabilization  of carbon metabolism mRNAs after repletion
of nitrogen following 16 hours of starvation \cite{Tesniere2017}. We do
not detect destabilization of \textit{PGK1} mRNA in our study and note that
the basal half-life of 6.2 minutes estimated in our study is similar
to the accelerated rate reported by \cite{Tesniere2017}. %inlineciteTesniere2017

%To identify the factors that underlie accelerated mRNA degradation, we
%developed a global \textit{trans}-factor screen using mRNA FISH, FACS, and
%sequencing. 
%BFF identified mutants in the Lsm1-7p/Pat1p
%complex as having elevated \textit{GAP1} mRNA levels both before and after the
%upshift.
%Given that the \textit{GAP1} mRNA is destabilized during this
%transition we suspect that these core mRNA degradation factors are
%directly involved. 
%Because factors associated with the Lsm1-7p/Pat1p
%complex are also involved in processing-body formation we looked for
%processing-body dynamics during the nitrogen upshift, but did not see
%qualitative changes in Dcp2-GFP distribution (raw data available in
%supplement). However, it has been proposed that pre-existing mRNPs
%seed the formation of processing-bodies \cite{Lui2014}, thus the
%phenotype may require assays at a finer spatial scale to eliminate
%this possibility \cite{Rao2017}. Interestingly, during
%cross-comparisons with a recent dataset exploring mRNA localization to
%RNP condensates \cite{Khong2017} we found that the set of
%destabilized transcripts in the label-chase experiment are on average
%longer in CDS and have an increased codon-optimality, two factors that
%were shown to be associated with differences in stress-granule
%localization of mRNA \cite{Khong2017}. 

%Regulated changes in mRNA stability can be mediated by RBP binding the
%3’ UTR of specific transcripts.  However, cis element analysis are
%inconsistent with a role for known RBPs in the observed
%destabilization. In particular, Puf3p motifs are de-enriched from the
%destabilized set. We failed to detect new 3’ UTR sequence motifs that
%are enriched in destabilized transcripts, but these 
Destabilized
transcripts are enriched for a binding motif of Hrp1p in
the 5’ UTR. This essential component of mRNA cleavage for
poly-adenylation in the nucleus has also been shown to shuttle to the
cytoplasm and bind to amino-acid metabolism mRNAs
\cite{Guisbert2005}. In addition, \textit{HRP1} has been shown to 
interact genetically to
mediate nonsense-mediated decay (NMD) of a \textit{PGK1} mRNA harboring a
premature stop-codon \cite{Gonzalez2000} or a \textit{cis}-element spanning
the 5’ UTR and first 92 coding bp of \textit{PPR1} mRNA \cite{Kebaara2003}.
A potential role for Hrp1p sites in the 5' UTR of destabilized mRNAs
warrants further investigation. 

Using BFF we identified mutants in the Lsm1-7p/Pat1p
complex as having elevated \textit{GAP1} mRNA levels both before and after the
upshift, consistent with their central role in mRNA 
metabolism. Out experiments using \textit{GAP1} normalized to
\textit{HTA1} demonstrate that increased mRNA abundance before the upshift is
likely a global effect (Fig~\ref{figure5}C). 
However, these mutants have a defect in clearance of \textit{GAP1}
mRNA upon the upshift, pointing to their function as mRNA clearance factors.
%and the assay demonstrates that associated decapping factors 
%\textit{EDC} and \textit{SCD6} have specific effects (Fig~\ref{figure5}D).
%Given that the \textit{GAP1} mRNA is destabilized during this
%transition we suspect that these mRNA degradation factors are
%directly involved. 
%While we found that the \textit{edc3}$\Delta$ mutant has defects in
%clearance of \textit{GAP1}, we also 
Additionally, we found that \textit{edc3}$\Delta$, 
\textit{scd6}$\Delta$,
%mutant shares a phenotype of reduced \textit{GAP1} mRNA
%expression during nitrogen limitation and reduced rate of \textit{GAP1} mRNA
%clearance with a 
and \textit{tif4632}$\Delta$ have impaired \textit{GAP1} mRNA clearance.
As these factors are involved in translational regulation they suggest
a role for this process in the differential stability of \textit{GAP1} mRNA
(Fig~\ref{figure5}D).
Deletion of the 5' UTR
of \textit{GAP1} also impairs its clearance (Fig~\ref{figure5}E). 
The 5' UTR deletions do not include the TATA box (ending at -179) or
GATAA sites (nearest at -237) responsible for NCR GATA-factor
regulation of \textit{GAP1} \cite{Stanbrough1996}.
Decreased steady-state expression of \textit{GAP1} mRNA in 5' UTR
mutants suggests that interactions of
factors with \textit{cis}-elements in the 5’ UTR might be responsible for
stabilizing \textit{GAP1} mRNA during nitrogen limitation, although the 
truncation of the 5' sequence may be enough to inhibit translation 
initiation by virtue of the shorter length \cite{Arribere2013}.
Elements in the 5’ UTR have
been demonstrated to modulate \textit{GAL1} mRNA stability
\cite{Baumgartner2011} and destabilize \textit{SDH2} mRNA upon glucose
addition, perhaps due to the competition between translation
initiation and decapping mechanisms \cite{De_la_Cruz2002}.
Interestingly, both \textit{GAP1} and \textit{SDH2} 
share the feature of a second start
codon downstream of the canonical start \cite{Neymotin2016} and
we have previously found that mutation of
the start codon of \textit{GAP1} results in lower
steady-state mRNA abundances \cite{Neymotin2016}.
This
%in light of recent analyses further highlighting the contribution of
%translation dynamics to mRNA stability 
%\cite{Presnyak2015,Neymotin2016,Cheng2017}, 
suggests a mechanism of degradation through dynamic changes in 
translation initiation that triggers decapping of \textit{GAP1} 
and other mRNA. 
%However, the deletion of \textit{SCD6} would be
%expected to promote translation of mRNA on the basis of it’s measured
%repressive activity in cell extracts, suggesting that if Scd6p does play a role
%that it may be specified by some condition-specific modulation of its
%activity \cite{Rajyaguru2012,Poornima2016}. 
Future work interrogating
this possible interaction of translational status and mRNA
stability during dynamic conditions could also expand our understanding of
the relationship between these two processes.

To our knowledge, this is the first time mRNA abundance has
been directly estimated using a SortSeq approach, although 
%sorting on indirect markers or 
using mRNA FISH and FACS to enrich subpopulations of cells has been
previously reported \cite{Klemm2014,Hanley2013,Sliva2016}. This
approach could be used with other barcoding mutagenesis technologies,
like transposon-based or Cas9 mediated perturbations, to
systematically test the genetic basis of transcript dynamics.
%phenotypes. A strategy combining this technology with transcriptomics
%as a high-dimensional marker could accelerate unbiased investigation
%of cellular signalling pathways \cite{Gapp2016}. Additionally, 
The use of branched-DNA mRNA FISH, or other methods
\cite{Rouhanifard2017}, allows for mRNA abundance estimation without
requiring genetic manipulation which makes it suitable for a variety
of applications. 
%While the cell wall
%of yeast makes optimization crucial to this assay, future development
%of hybridization protocols may improve accuracy and make the assay
%more robust \cite{Richter2017,Wadsworth2017}. 
Furthermore, our methods for library construction from limited material
should permit accurate
quantification of pooled barcode libraries with small inputs, 
expanding the possibilities for flow cytometry markers to fixed-cell assays.

Why is \textit{GAP1} subject to multiple layers of repression upon
a nitrogen upshift, 
including at the level of transcript synthesis, transcript degradation,
protein maturation, and post-translational inactivation? Given the
strong fitness cost of inappropriate activity \cite{Risinger2006},
this overlap could ensure mechanistic redundancy for robust repression in
the face of phenotypic or genotypic variation. Alternatively, it could
reflect a systematic need to free ribonucleotides or
translational capacity, or some as yet uncharacterized
process.
%, or could simply be an effect of some unrelated function. 
%While this question remains open,
%we have made progress towards this goal by identifying 
%factors required for its accelerated degradation.
%of decapping associated with the Lsm1-7p/Pat1p complex play a role. 
Future work aimed at determining the adaptive basis of accelerated
mRNA degradation will serve to illuminate the functional role of
post-transcriptional gene expression regulation.
%dissecting this
%mechanism and contrasting the dynamic process of mRNA destabilization
%during other growth transitions would greatly inform our understanding
%of mRNA stability specification at steady-state, possibly in light of
%the relationship between translation and stability of mRNAs. 

\section*{Materials and methods}

\subsection*{Availability of data and analysis scripts}

\label{codeanddata}

Computer scripts used for all analyses are available as a git repository
on GitHub
(\texttt{\url{https://github.com/darachm/millerBrandtGresham2018}}),
sequencing fastq files are available in the Short Read Archive 
(SRP142613),
and all data along with intermediate analysis files are available in 
zip archives archived in the Open Science Framework 
(\texttt{\url{https://osf.io/7ybsh/}}).
With these archives, all analyses downstream of alignments
can be re-run completely using the provided GNU Makefile.
A Shiny application is also available to explore the two main 
datasets in this paper at
\url{http://shiny.bio.nyu.edu/users/dhm267/}. It
is also included in the OSF repository as a separate zipped 
archive for local installation and long-term archiving. 
Consult the \texttt{README.md} file in the git repository for more 
specific instructions, and the \texttt{html\_reports.zip} archive
for Rmarkdown-generated reports of all R analyses used.

\subsection*{Media and upshifts of media}

The nitrogen-limited media used is a minimal media
supplemented with various salts, metals, minerals, vitamins, and
2\% glucose, as previously described \cite{Airoldi2016,Brauer2008}. 
For proline limitation, 
the media was made with 800μM L-proline as the sole
nitrogen source.
YPD media was made using standard recipes \cite{Amberg2005}.
All growth was at 30$^{\circ}$C, in an air-incubated 200rpm shaker  
using baffled flasks with foil caps, or roller drums for 
overnight cultures in test tubes.
For glutamine upshift experiments, 
400μM L-glutamine was added from a 100mM stock solution dissolved 
in MilliQ double-deionized water and filter sterilized.
All upshift experiments were
performed at a cell density of 1-5 million cells per mL,
in nitrogen-limited media in which untreated cultures 
saturate at approximately 30 million cells per mL. 
For all experiments, 
a colony was picked from a YPD plate and grown in a 5mL
nitrogen-limited (proline)
pre-culture overnight at 30$^{\circ}$C, then innoculated into
the experimental culture from mid-exponential phase.

\subsection*{Strains}

See \nameref{strains} for details.  
The wild-type strain used is FY4, a S288C derivative. 
The pooled deletion collection is as published
\cite{Vandersluis2014}. %inlineciteVandersluis2014
For all experiments with single strains, colonies were struck 
from a -80$^{\circ}$C frozen stock onto YPD (or YPD+G418 for
deletion strains) to isolate single colonies.
For pooled experiments we inoculated directly into nitrogen-limited
(proline) media
from aliquots of frozen glycerol stocks.

Strains with deletions 5' of the start codon and 3' of the stop
codon were generated by the "delitto-perfetto" 
method \cite{Storici2006}, 
by inserting the pCORE-Kp53 cassette
at either the 5' or 3' end of the coding sequence, then transforming
with a short oligo product spanning the deletion junction and
counter-selecting against the cassette using Gal induction of p53.
Strains were generated and confirmed by Sanger sequencing,
and traces are available in directory \texttt{data/qPCRfollowup/} 
within the data zip archive (\nameref{codeanddata}).

\subsection*{Measurement of growth during upshift}

A single colony of FY4 was inoculated in 5mL nitrogen-limited
(proline)
media and grown to exponential phase, then back diluted in
nitrogen-limited (proline) media
in a baffled flask. 
Samples were collected into an eppendorf, sonicated,
diluted in isoton solution, and analyzed with a Coulter Counter Z2
(Beckman Coulter).
%within approximately 5 minutes 
%of collection from flask. Output files were read and processed by a 
%script (in the git repository) to yield counts per diameter bin, 
%and minimum size particle 
%gating was used with blank controls to exclude measurement noise.
%<!--We are only presenting one replicate, as similar results  
%have been obtained many times by our lab and 
%we chose to generate a rapidly sampled time course to best represent 
%the upshift. Seriously, people have seen this since Slator 1918.
%-->

\subsection*{Re-analysis of microarray data} 

%Supplemental files from \cite{Airoldi2016} were  %inlineciteAiroldi2016
%downloaded, read
%into \texttt{localc} (an open-source spreadsheet software), 
%a small Excel-generated auto-correction error was 
%fixed ("Oct-1" -> "OCT1"), and the file saved as a CSV. Microarray 
%intensity ratios were processed with 
Gene expression data \cite{Airoldi2016} were analyzed using
\texttt{pcaMethods} \cite{Stacklies2007} to perform a SVD PCA on scaled data. 

\subsection*{qPCR}

Samples were collected before, during the first ten minutes of
the nitrogen upshift (Fig~\ref{figure3}),
or at ten minutes after the upshift (Fig~\ref{figure5}).
For the experiments described in Fig~\ref{figure5}, all work
was done in biological replicates.
Each 10mL sample was collected by vacuum onto a 25mm nylon filter
and frozen in liquid nitrogen.
RNA was extracted by adding 400μL of TES buffer
(10mM Tris (7.5pH), 10mM EDTA, 0.5\% SDS)
and 400μL of acid-phenol, vortexing vigorously and incubating at 
65$^{\circ}$C for an hour with vortexing every 20 minutes. 
For Fig~\ref{figure3} only, at the beginning of this extraction incubation
we added 10μL of a 0.1ng/μL in-vitro synthesized spike-in 
mRNA BAC1200 (as generated
for the label-chase RNAseq (\nameref{writeup2}), 
but without 4-thiouridine). 
All samples were separated by centrifugation and extracted again 
with chloroform on a 2mL phase-lock gel tube (5Prime \#2302830). 
After ethanol precipitation of the aqueous layer, 
RNA was treated DNAse RQ1 (Promega M610A) according to manufacturer
instructions, then the reaction heat-killed at 65$^{\circ}$C for 
10 minutes after adding a mix of 1:1 0.5M EDTA and RQ1 stop-solution.
The resulting RNA was
cleaned with a phenol-chloroform extraction and ethanol 
precipitated.
All samples were hybridized with RT primers by incubating the mixture at 80$^{\circ}$C for 
5 minutes then on ice for 5 minutes.
For Fig~\ref{figure3} 2μg RNA was primed with 2.08ng/μL
random hexamers (Invitrogen 51709) and 
2.5mM total dNTPs (Promega U1511),
while for Fig~\ref{figure5} 1μg RNA was primed with 
5.6mM Oligo(dT)18 primers (Fermentas FERSO132) and
0.56mM total dNTPs (Promega U1511).
These mixtures were combined with 1/10th 10x M-MulvRT buffer (NEB M0253L), 
1/20th volume RNAse-OUT (Invitrogen 51535), and 1/20th volume M-MulvRT (NEB M0253L). 
A negative control with no reverse-transcriptase enzyme was also prepared
and analyzed in the qPCR reaction.
The reaction proceeded for 1 hour at 42$^{\circ}$C, 
then was heat-killed at 90$^{\circ}$C
before diluting 1/8 with hyclone water (GE SH30538). 
This dilution was used as direct
template in 10μL reactions with SybrGreen I Roche qPCR master-mix
(Roche 04 707 516 001) for measurement on a Roche Lightcycler 480. 
For Fig~\ref{figure3}, we used primers 
DGO230 and DGO232 to quantify \textit{GAP1} and 
primers DGO605 and DGO606 to quantify the synthetic spike-in BAC1200.
For Fig~\ref{figure5}, we used primers
DGO229 and DGO231 to quantify \textit{GAP1} and
primers DGO233 and DGO236 to quantify \textit{HTA1}.  
See \nameref{primers} for sequences.
These were run on a Roche480 Lightcycler, 
with a max-second derivative estimate
of the cycles-threshold (the $C_p$ value output by analysis) used 
for analysis by scripts included in the git repo 
(\nameref{codeanddata}).
Linear regression of the log-transformed values was used to quantify
the dynamics and assess significance of changes in expression
levels or rates of change.

\subsection*{Microscopy}

Cells hybridized with \textit{GAP1} mRNA FISH Affymetrix probes (as 
described in detail in \nameref{writeup4}) were sorted 
with a BD FACSAria II based on emission area from a
660/20nm filter with a 633nm laser activation
into four gates for the 3-minute post-shift timepoint. 
These were sorted using PBS sheath fluid at room-temperature into
poly-propylene FACS tubes, vortexed and applied to
poly-L-lysine-treated coverslips. 
Images were acquired on a DeltaVision scope, with FISH fluorescence 
detected in the "Cy5" channel (632/22nm excitation, 676/34nm emission)
and the "Visible" light collected as bright-field illumination 
captured with a polarized objective. Raw images available in the
"microscopy" zip archive (\nameref{codeanddata}).

\subsection*{Microscopy of Dcp2-GFP}

To look for processing-body dynamics in response to
a nitrogen upshift, we used strain DGY525, which is FY3
containing plasmid pRP1315 (gift from Roy Parker).
Samples were collected before and following a nitrogen upshift
(4, 10, 12, 19, or 25 minutes),
from exponential growth in YPD, or 10 minutes after resuspending
YPD-grown cells in DI water.
All samples were collected by centrifugation at 10,000g for 30 seconds, 
aspirating most supernatant, then centrifugation for 20 seconds
and aspirating all media. Each pellet was 
immediately resuspended in 4\% PFA 
(diluted from EMS 16\% PFA ampule RT15710) 
with 1x PBS ( NaCl 8g/L, KCl 0.2g/L, Na$_2$HPO$_4$ 1.42g/L, 
KH$_2$PO$_4$ 0.24g/L) 15 minutes on bench, 
then spun at 10,000g for 1 minute, aspirated, 
then washed once and resuspended with 1x PBS. 
Samples were kept on ice, then put onto a coverslip
for imaging on a DeltaVision scope. Raw images available in the
microscopy zip archive (\nameref{codeanddata}).

\subsection*{4tU label-chase and RNA sequencing}

The methods and analysis are detailed in \nameref{writeup2},
including protocols and manufacturer information,
and all data and code are available as described
(\nameref{codeanddata}).

FY4 was grown in nitrogen-limitation conditions overnight 
with a 50μM:50μM mixture of 4-thiouracil:uracil.
This culture was split, then 4mM uracil was added to chase the
4-thiouracil label (a 41-fold excess of uracil).
30mL samples of the culture were taken by filtration onto 25mm
nylon filters and flash-frozen in eppendorfs. 
After letting the chase proceed for 12.5 minutes, 
we added glutamine from 
a 100mM stock (dissolved in water) to a final concentration 
of 400μM to one flask, or an equal amount of water to the 
control flask. 
Samples were extracted using a hot acid-phenol method,
with equal volume of synthetic spike-ins  added to each RNA 
extraction reaction.
4tU-containing spike-ins (polyadenylated coding sequences from
\textit{B. subtilus} and \textit{C. elegans}) were synthesized
\textit{in-vitro} as previously described \cite{Neymotin2014}. 
RNA was reacted with 
MTSEA-biotin to conjugate biotin to the 4-thiouracil-containing
RNA, then purified using streptavidin beads.
Fractionated RNA was depleted of rRNA using a RiboZero kit.
RNA samples were converted into Illumina sequencing libraries using
a strand-specific (UNG) protocol, ligating adapters that contain
UMI's \cite{Hong2017}.
Libraries were pooled and sequenced by the NYU Genomics
Core sequencing facility on an Illumina NextSeq.
Following base-calling and sample demultiplexing by NYU GenCore,
the sequencing reads were trimmed using 
\texttt{cutadapt} \cite{Martin2011}
aligned using \texttt{tophat2} \cite{Kim2013} to a reference genome
that included the yeast reference genome (assembly R64-2-1) and
spike-ins, filtered for mapping-quality and length using
\texttt{samtools} \cite{Li2009}, deduplicated 
with \texttt{umi\_tools} \cite{Smith2017}
and feature counting was performed using 
\texttt{htseq-count} \cite{Anders2015}.
Feature counts for yeast mRNAs were normalized to synthetic spike-ins, 
using the fitted values from a log-linear model of spike-in abundance
increase (see Results, and \nameref{writeup2}).
The rate of mRNA degradation and changes in this rate
was quantified assuming an exponential model 
(\nameref{writeup2}) and fit as a linear model to log 
transformed data.
Significant changes in mRNA degradation rates were defined using 
a FDR \cite{Storey2015} less than 0.01 and a doubling in degradation
rate (based on modeling detailed in \nameref{writeup2}).

\subsection*{Label-chase RNA sequencing \textit{cis} element analysis}

To detect if \textit{de novo} or known \textit{cis} elements were 
associated with destabilization upon a nitrogen upshift,
we used 
DECOD \cite{Huggins2011}, FIRE \cite{Elemento2007},
TEISER \cite{Goodarzi2012}, and the \#ATS pipeline \cite{Li2010}. 
We also scanned for association with
RBP binding sites from the CISBP-RNA database
\cite{Ray2013} using AME from the MEME suite
\cite{McLeay2010}. 
Final plots in the supplement were made using motif scans with 
GRanges \cite{Lawrence2013}.
Analysis was done using coding sequence and 
four different definitions of untranslated regions 
(200bp upstream of the start codon or downstream of the stop codon, 
the largest detected isoform in TIF-seq data \cite{Pelechano2014},
or the most distal detected gPAR-CliP sites in exponential-phase 
or nitrogen-limited growth \cite{Freeberg2013}).

\subsection*{Barseq after FACS after mRNA FISH (BFF)}

The methods and analysis are detailed in \nameref{writeup4},
including motivation, protocols, and manufacturer information.

An aliquot of the prototrophic deletion collection
\cite{Vandersluis2014} was thawed and diluted, with 
approximately 78 million cells added to 500mL of nitrogen-limited
(proline) media 
in a 1L baffled flask. The culture was shaken at 30$^{\circ}$C overnight, 
then split into three flasks (A, B, and C). 
After three hours (at mid-exponential)
we collected samples of 30mL culture filtered onto a 25mm filter and
flash-frozen in an eppendorf in liquid nitrogen. 
We sampled in steady-state growth (pre-upshift) and  
10.5 minutes after adding 400μM glutamine (post-upshift).
Samples of the pool were fixed with formaldehyde
(4\% PFA diluted in PBS from 10mL aliquot, 
buffered, 2 hours room-temperature) and digested
with lyticase (in BufferB with VRC 37$^{\circ}$ 1 hour), 
\cite{Mcisaac2013}, and permeabilized with ethanol at
4$^{\circ}$ overnight.
Samples were processed with a Affymetrix Quantigene Flow RNA kit 
(product code 15710) designed to target
\textit{GAP1} mRNA and labelled with Alexa 647.
The hybridization was performed using a modified
version of the manufacturer's protocol (Appendix
\nameref{writeup4}, including a DAPI staining step.
Samples were sonicated, then run through a BD FACSAria II.
Cells were gated for singlets and DAPI content 
(estimated 1N or more), then sorted based on signal detected with a
%emission area from a
660/20nm filter with a 633nm laser activation
into four gates within each timepoint. %, across replicates.
Cells were sorted using PBS sheath fluid at room-temperature, into
poly-propylene FACS tubes, then stored at -20$^{\circ}$C.
For each gate, cells were collected via centrifugation and genomic
DNA extracted by NaCl reverse-crosslinking at 65$^{\circ}$C,
inspired by \cite{Klemm2014}, with %inlineciteKlemm2014
subsequent proteinase K and RNase A digestions.
Genomic DNA was split into three reactions and assayed using a 
modified barseq protocol (\nameref{writeup4}).
See the supplementary write-up \nameref{writeup4} for 
detailed protocols and rationale. %, and a discussion of dimers.
Barseq libraries were submitted to the NYU Genomics Core 
for sequencing using a 1x75bp run on a Illumina NextSeq.

\subsection*{Analysis of BFF sequencing results}

We devised a pipeline to quantify barcodes using the UMI sequence 
incorporated in the first round of
amplicon priming, and benchmarked on \textit{in silico}
simulated datasets (\nameref{writeup4}).
Briefly, raw FASTQ files are processed with SLAPCHOP
(\url{https://github.com/darachm/slapchop}) 
which uses pair-wise alignment
\cite{Cock2009} to filter, extract UMIs from 
variable positions, and extract barcodes. % into different fields.
We demultiplexed using a perl script, and aligned 
strain barcodes to a reference barcode index \cite{Smith2009} 
using \texttt{bwa mem} \cite{Li2013}. Barcodes were counted
and then we used UMIs with the label-collision correction of 
\cite{Fu2011} to quantify the proportion of each mutant in the %inlineciteFu2011
sample. These relative counts and the FACS data
(the sorted events per bin) were used to estimate the distribution 
of each mutant across the four gates in each timepoint.
We filtered for strains detected in at least three bins,
and fit a log-normal distribution using \texttt{mle} in R
\cite{Team2000}. The mean of this distribution %inlineciteTeam2000
was used as the expression value of \textit{GAP1} mRNA in plots and
GSEA analysis using \texttt{clusterProfiler} \cite{Yu2012}.

\clearpage

\section*{Figure captions}

\begin{figure}[!ht]
  \caption{
%\includegraphics[width=\textwidth]{output/Figure1.png}
  \textbf{Dynamics of physiology and transcriptome remodeling
  during a nitrogen upshift.}
  \textbf{A)} 
  400μM glutamine was added to a
  culture of yeast cells growing in minimal media containing 800μM
  proline as a sole nitrogen source. Measurements
  of culture density across the upshift are plotted. 
  Dotted lines denote linear regression of the
  natural log of cell density against time before the upshift and 
  after the 2 hour lag. \textbf{B)} Average cell size during the
  same experiment.
  Dotted lines denote the mean cell diameter before the upshift
  and after the 2 hour lag. 
  \textbf{C)} PCA analysis of global
  mRNA expression in steady-state chemostats and following an upshift
  \cite{Airoldi2016}. Steady-state nitrogen-limited chemostat
  cultures maintained at different growth rates (colored circles)
  primarily vary along principal component 2. Expression following a
  nitrogen-upshift in either a chemostat (squares) or batch culture
  (triangles) show similar trajectories and vary along
  both components. Grey lines depict the major 
  trajectory of variation for steady-state and upshift experiments.
  \label{figure1}
  }
\end{figure}

\begin{figure}
%  \includegraphics[width=\textwidth]{output/Figure2.png}
\caption{  \textbf{Global mRNA stability changes following a nitrogen upshift.}
  \textbf{A)} 4tU-labeled mRNA from each gene was measured over time, before and
  after the addition (vertical dotted line) of glutamine 
  (nitrogen-upshift) or water (mock). Linear regression models were 
  fit to the data with a rate before the upshift (solid line) 
  and a change in rate after glutamine addition (dashed line). 
  \textit{HTA1} is not significantly destabilized, 
  whereas mRNAs encoding NCR-regulated transporters or 
  pyruvate and trehalose metabolism components are 
  significantly destabilized. 
  Plots for all genes are available in the associated Shiny application
  (Methods).
  \textbf{B)} Comparison between the pre-upshift mRNA
  degradation rate (y-axis) and the post-upshift mRNA degradation rate
  (x-axis). % Positive values result from noise on the slope estimate.
  Details of modeling are in \nameref{writeup2}.
  \textbf{C)} Comparison between changes in mRNA expression following
  upshift \cite{Airoldi2016} (y-axis) and the post-upshift
  degradation rate (x-axis). % Both plots share the same x-axis.
  Transcripts that are significantly destabilized are colored red, and
  shown with red rug-marks in the marginal histograms.
  \label{figure2}
  }
\end{figure}

\begin{figure}
%\includegraphics[width=\textwidth]{output/Figure3.png}
\caption{
  \textbf{\textit{GAP1} mRNA dynamics measured by flow cytometry.}
  \textbf{A)} 
  GAP1 mRNA following upshift measured using RT-qPCR, relative
  to an external spike-in mRNA standard. The dashed line is fit
  to points 2 minutes after the upshift. 
  \textbf{B)} Flow
  cytometry of wild-type yeast probed for \textit{GAP1} mRNA
  in nitrogen-limited conditions and
  following an upshift. The vertical grey lines indicate FACS gate
  boundaries used for cell sorting. 
  \textbf{C)} Representative cells from each bin sorted from the
  experiment in panel B. 
  \textbf{D)} Quantification of microscopy data. 
  Each black dot represents a single cell. 
  The mean number of foci per cell in each bin from panel B
  is displayed as a red point.
  \label{figure3}
  }
\end{figure}

\begin{figure}
%\includegraphics[width=\linewidth]{output/Figure4.png}
\caption{
  \textbf{BFF estimates of \textit{GAP1} mRNA abundance per mutant.}
  \textbf{A)} Flow cytometry analysis of \textit{GAP1} mRNA 
  abundance in the prototrophic
  deletion collection (n=3,230 mutants) before and after the upshift. 
  The vertical gray
  lines denote boundaries of the four FACS gates. 
  Biological replicates are indicated by color. 
  \textbf{B)} Measurements for individual genes before and
  after the upshift. Pseudo-events per strain per bin 
  are on the y-axis.
  Black dashed lines indicate maximum-likelihood 
  fits of a log-normal to pseudo-events within each bin 
  for each mutant. For plotting purposes, points are positioned on the
  x-axis at the average signal for the library in that bin.
  Colors are as in panel A. 
  \textbf{C)} Distribution of modeled mean GAP1 mRNA levels
  for each mutant.
  \textbf{D)} The mean \textit{GAP1} mRNA expression levels 
  fit using all replicate data for individual mutants before and after
  the upshift are shown as points connected by a line, colored
  according to the type of gene. 
  The background violin plot shows the distribution 
  of all 3,230 mutants.
  Plots for all mutants are available in the associated Shiny application
  (Methods).
  \label{figure4}
  }
\end{figure}

\begin{figure}
%\includegraphics[width=\textwidth]{output/Figure5.png}
\caption{
  \textbf{Disrupting the Lsm1-7p/Pat1p complex and translational
  regulation impairs clearance of \textit{GAP1} mRNA.}
  \textbf{A)} In the background is the distribution of 
  fit \textit{GAP1} mRNA mean expression levels for all mutants
  in the pool. Indicated by colored points and lines are the means for
  individual knockout strains, as labeled.
  \textbf{B-E)}, \textit{GAP1} mRNA relative to
  \textit{HTA1} mRNA before and 10 minutes after a glutamine upshift, 
  in biological replicates. Lines are a log-linear regression fit. 
  Points are dodged horizontally for clarity, but timepoints for
  modeling and for drawn lines are 0 and 10 minutes exactly.
  Wild-type is FY4, and each estimate of the \textit{GAP1}/\textit{HTA1}
  ratio is normalized to the average ratio measured of FY4 at t=0 
  for that qPCR batch.
  \textbf{B)} \textit{xrn1}$\Delta$, \textit{ccr4}$\Delta$,
  \textit{pop2}$\Delta$ are all defective in \textit{GAP1} mRNA clearance (p-values $<$ 0.004).
  \textbf{C)} \textit{lsm1}$\Delta$ and \textit{lsm6}$\Delta$ are 
  slowed in \textit{GAP1} mRNA clearance (p-values $<$ 0.0132 and 0.0299, respectively).
  \textbf{D)} \textit{edc3}$\Delta$ is slowed in \textit{GAP1} mRNA clearance 
  (p-value $<$ $10^{-4}$).
  \textit{scd6}$\Delta$ and \textit{tif4632}$\Delta$ are slowed in
  \textit{GAP1} mRNA clearance (p-values $<$ $10^{-5}$) and have lower levels of expression
  before the upshift (p-values $<$ 0.003).
  \textbf{E)} A deletion of 150bp 3' of \textit{GAP1} stop codon has
  no significant effect, but a deletion of 100bp 5' of the start
  codon has a defect in \textit{GAP1} mRNA clearance 
  (p-value $<$ $10^{-4}$) and lower level of expression before the upshift
  (p-value $<$ 0.0015).
  \label{figure5}
  }
\end{figure}

\clearpage

\section*{Supporting information}

% Include only the SI item label in the paragraph heading. Use the \nameref{label} command to cite SI items in the text.

%
%
%

\paragraph*{S1 Appendix.}
\label{writeup2}
\textbf{Supplementary file with experimental rationale, details,
and protocol for the label-chase experiment.}


\paragraph*{S2 Appendix.}
\label{writeup4}
\textbf{Supplementary file with experimental rationale, details, 
and protocol for the BFF experiment.}


%
%
%

\paragraph*{S1 Table.} 
\label{microarrayPCAloadings}. 
\textbf{Principal components loadings of the microarray samples reprocessed from Airoldi et. al. 2016.}

\paragraph*{S2 Table.} 
\label{microarrayPCAgsea}. 
\textbf{Gene set enrichment analysis of loadings on principal components one and two.}

\paragraph*{S3 Table.}
\label{dme211raw}
\textbf{Raw counts of labeled mRNA quantified by RNAseq in label-chase experiment.}

\paragraph*{S4 Table.}
\label{dme211filterDirect}
\textbf{Filtered label-chase RNAseq data for modeling, normalized directly within sample.}

\paragraph*{S5 Table.}
\label{dme211filterModel}
\textbf{Filtered label-chase RNAseq data for modeling, normalized by modeling across samples.}

\paragraph*{S6 Table.}
\label{dme211resultsModel}
\textbf{Degradation rate modeling results, from data normalized across samples.}

\paragraph*{S7 Table.}
\label{dme211goAndKegg}
\textbf{Enriched GO and KEGG terms within the set of mRNA destabilized upon a nitrogen upshift, across sample normalization.}

\paragraph*{S8 Table.}
\label{dme209rawCountsGates}
\textbf{Raw counts of events in each bin
in the BFF experiment, and gate settings for the observations.}

\paragraph*{S9 Table.}
\label{dme209pooledFits}
\textbf{Parameters of all 3230 models fit to filtered data.}

\paragraph*{S10 Table.}
\label{dme209gsea}
\textbf{Gene-set enrichment analysis results using \textit{GAP1} estimates.}

\paragraph*{S11 Table.}
\label{strains}
{\bf Strains used in this study.} 

\paragraph*{S12 Table.}
\label{primers}
{\bf Primers used in this study.} 
All primers were synthesized by Integrated DNA Technologies (IDT).

Barseq multiplexing barcode sequences and index numbers available in 
the file \texttt{data/dme209/sampleBarcodesRobinson2014.txt} within 
the data zip archive (\nameref{codeanddata}).



%
%
%

\paragraph*{S1 Fig.}
\label{longTermPCA}
\textbf{The long-term transcriptome dynamics of a glutamine
upshift.}
Principal components analysis (SVD) of microarray data from 
\cite{Airoldi2016}.  %inlineciteAiroldi2016
Colored points are from steady-state chemostats grown in
limitation for various nitrogen sources, at different growth rates.
Time-series experiments are show in grey points, connected by lines,
and line-type is the type of upshift (in batch or in chemostat).

%

\paragraph*{S2 Fig.}
\label{timeplot}
\textbf{The timing of changes in the component loadings of microarray samples.}
The changes in loadings in the first two principal 
components of microarray data analyzed from \cite{Airoldi2016},
as in \textbf{Fig 1C}, for chemostat upshift experiments only.
In the chemostat, addition of 400μM glutamine has a more 
pronounced response than 40μM glutamine addition, 
but all have a sharp response in short timescale.

%

\paragraph*{S3 Fig.}
\label{timeplot}
\textbf{Most of the variation is explained in the first two principal components.}
\textbf{A)} From the PCA analysis used in \textbf{Fig 1C}, 
the variance explained
is plotted for the principal components, showing a steep 
decrease in explained variance after the first two components.
\textbf{B)} The third and fourth components explain a small amount
of the variance and are not readily interpretable with respect to
the transcriptome changes in response to the changing growth rates
in either steady-state or dynamic conditions.

%

\paragraph*{S4 Fig.}
\label{kkdComparison}
\textbf{Comparison between rates of mRNA abundance changes and 
  mRNA stability.}
Comparisons of measured mRNA degradation from this study with mRNA abundance change rates
from \cite{Airoldi2016}. 
Pre-upshift degradation rates (top) don't explain the 
abundance change. 
The degradation rate changes (middle, difference between pre
and post upshift) and the post-upshift rates (bottom) are anti-correlated
with the abundance changes.


\paragraph*{S5 Fig.}
\label{comparisonESR}
\textbf{Many of the destabilized mRNA are members of the ESR-up regulon.}
Comparisons of degradation rates from this study with mRNA abundance change rates
from \cite{Airoldi2016}. Destabilized transcripts are colored based on %inlineciteAiroldi2016
their membership in the ESR gene set, as described in the supplement 
of \cite{Brauer2008}.  %inlineciteBrauer2008
Many of the destabilized set are "ESR-up" genes, as they
are increase in expression in response to stresses.


\paragraph*{S6 Fig.}
\label{comparisonDestabilized}
\textbf{Scatter plot of significantly destabilized transcripts.}
For each transcript the x-axis is
the rate of degradation rate post-upshift and the y-axis is the mRNA abundance
change rate \cite{Airoldi2016} after the upshift.
%These values were modeled to normalized
%sequencing signal (x-axis) and normalized microarray ratio (y-axis). 
The dashed line is a 1:1 line of equality. 

\paragraph*{S7 Fig.}
\label{compareSix}
\textbf{Examples of individual mRNA with counteracting changes in 
mRNA synthesis and degradation.}
For several transcripts we found an increased rate of degradation
post-upshift (red) compared to before the upshift (blue) but minimal
changes in abundance (black).
Each dataset is normalized to intersect at the same t=0 intercept.
%examples of the slowest decreasing (in the microarray fits)
%transcripts, we plot the microarray (abundance) and sequencing (decaying labeled
%abundance) data normalized to be on the same relative y-axis scale (subtracted
%t\_0 y-intercepts of fits).
%Black is the microarrays, so abundance. Blue is the
%water-upshifted samples. Red is the glutamine-upshifted samples. This shows how
%transcripts can be destabilized (blue slope changing to red slope) without being
%greatly repressed (black line is not of same slope as red line), showing that
%synthesis and degradation can contradict each other for some transcripts in this system.
%Black is the microarrays, so total abundance. Blue is the 4tU-labeled
%signal in water-upshifted samples. 
%Red is the 4tU-labeled signal in glutamine-upshifted samples. 
%Destabilization does not necessarily result in a rapid clearance
%of the mRNA.

\paragraph*{S8 Fig.}
\label{lengthAndCodons}
\textbf{The destabilized set of transcripts is longer and has a higher frequency
of optimal codons than the rest of the transcriptome.}
Comparisons of destabilized mRNAs with the rest of the transcriptome.
\textbf{A)} 
Destabilized transcripts tend to have longer CDS lengths (p-value $<\;2\times10^{-5}$
by Wilcoxon rank sum test). 
\textbf{B)} 
On average, the destabilized transcripts have more optimal codons
than the rest of the transcriptome (p-value $<\;2\times10^{-8}$ Wilcoxon rank
sum test).
The fraction of optimal codons per feature
was obtained from the supplement of \cite{Khong2017} using definitions 
from \cite{Presnyak2015}.  


\paragraph*{S9 Fig.}
\label{hrp1}
\textbf{Enrichment of Hrp1p motif in 5' UTRs of destabilized
transcripts.}
Sequences were analyzed for RBP binding motif enrichment
using the AME program in MEME and significant hits confirmed using a
logistic model predicting destabilization based on motif content per sequence
length. Hrp1p is significantly ( p $<$ 0.0001 ) enriched in
the 5' UTRs of destabilized transcripts. Motif matches were
counted using the GRanges package for the 5' UTRs, 3' UTRs,
and coding sequence of transcripts using the largest isoforms detected in
\cite{Pelechano2014}. %inlinecitePelechano2014

%

\paragraph*{S10 Fig.}
\label{gap1Delete}
\textbf{\textit{GAP1} delete or omission of the targeting probe results in no \textit{GAP1} FISH signal.}
%\textit{GAP1} delete or omission of the targeting probe removes signal of GAP1 FISH.
Wild-type or \textit{GAP1}$\Delta$ cells were grown in proline-media.
%, which induces expression of \textit{GAP1}. 
As seen in the positive control there is heterogeneity in the
signal. 
%This is likely due to technical issues, namely fixation time. 

%

\paragraph*{S11 Fig.}
\label{pca}
\textbf{Principal components analysis of the abundance estimates using BFF.}
Each color is a type of sample, from low to high gates (with black denoting the
input samples before sort). Technical replicates are connected by dashed lines,
biological replicates are denoted with letters A B or C. 
The first two principal components
show the separation of gates by signal intensity and reflects that the lower
gates on the upshifted samples were very close (blue and red samples on far right
panel), within the distribution of the negative population. This is consistent
with their tight sampling of the "GAP1-off" population, as seen in
\textbf{Fig~\ref{figure4}A}.


\paragraph*{S12 Fig.}
\label{rarefaction}
\textbf{Rarefaction curve of UMI saturation.}
The solid-line curve denotes the theoretical expectation of total 
observations per UMI in a sample (x-axis) and the number of 
unique UMIs (y-axis). This curve shows how 
UMI-collisions are expected to depress the number of unique UMIs. Each 
point is from real data, with these
two numbers tabulated for each combination of a sample and strain barcode. We
see that these largely follow the curve of saturation of UMI-collisions, but
that it falls well below the expectation of independent UMI-collision, thus
we believe that there is an additional contribution of PCR-amplification noise
(PCR duplicates). 

%

\paragraph*{S13 Fig.}
\label{gluco}
\textbf{Knock-out mutants of negative regulators of 
gluconeogenesis are associated with higher \textit{GAP1} mRNA expression 
after the upshift.}
Knock-out mutants of negative regulators of gluconeogenesis are associated with
higher estimated \textit{GAP1} mean after the upshift, by GSEA analysis
of GO-terms (p-value $<$ 0.05).


\paragraph*{S14 Fig.}
\label{sulfate}
\textbf{Knock-out mutants of genes involved in sulfate assimilation
are associated with higher estimated \textit{GAP1} mRNA expression mean before the upshift.}
Knock-out mutants of involved in sulfate assimilation are associated with
higher estimated \textit{GAP1} mean before the upshift, by GSEA analysis
of GO-terms (p-value $<$ 0.05).


\paragraph*{S15 Fig.}
\label{prePredictPost}
\textbf{The relationship between mean \textit{GAP1} mean expression before the 
shift and after the upshift.}
%The relationship between the estimated mean before the upshift and after the
%upshift. 
Scatter plot of the estimated means, with marginal histograms along
top and right. Red vertical line on top histogram is a cut-off of
\textit{GAP1} mRNA induction for this analysis,
and is the mean of the fit to wild-type minus one standard deviation of that
distribution. The red linear regression line is fit to all points above this
threshold, in which expression was detected before the upshift.



\paragraph*{S16 Fig.}
\label{tco89}
\textbf{\textit{tco89}$\Delta$ and \textit{xrn1}$\Delta$ show
defects in \textit{GAP1} mRNA regulation in the BFF assay.}
%Data and fits for several mutants. 
\textit{xrn1}$\Delta$ mutant (left) is lowly abundant in
the library and is only observed in the highest bin of \textit{GAP1} signal, consistent
with the role of Xrn1p as a global exonuclease. 
\textit{tco89}$\Delta$ is the only detected member that would abrogate TORC1 activity.
This mutant (right) has elevated \textit{GAP1} mRNA before and after the upshift,
consistent with the role of TORC1 in repressing the NCR regulon. 

%

\paragraph*{S17 Fig.}
\label{bothutr}
\textbf{Two independent deletions of the \textit{GAP1} 5' UTR 
show the same phenotype.}
%During strain construction, 
A deletion of 152bp 5' of the start
codon was also generated. We tested \textit{GAP1} dynamics in this
strain as well, and found that it shares the same phenotype as a
100bp 5' UTR deletion. Methods are the same as in \textbf{Fig~\ref{figure5}E},
both 5' UTR deletes are slowed in clearance, ANCOVA p $<$ 0.05 .

%

\paragraph*{S18 Fig.}
\label{pbodyScope}
\textbf{Processing-body dynamics are not associated with the
nitrogen upshift.}
A strain containing a copy of Dcp2p-GFP expressed from a plasmid
was grown in conditions of exponential phase in YPD or 10 minutes of
starvation in water (first row) to confirm detection of processing-body
foci using Dcp2-GFP.
%result in the formation of processing-body foci of Dcp2-GFP.
We do not see either formation or dissolution of Dcp2-GFP foci during
the nitrogen upshift (bottom row).

\section*{Acknowledgments}

We would like to thank past and present members of the Gresham and 
Vogel labs for discussions and feedback, 
Andreas Hochwagen and Viji Subramanian 
for microscope access and training, Ken Birnbaum for helpful 
conversations and equipment usage, Evelina Tutucci for demonstrating 
mRNA FISH, Andres Mansisidor and Matt Robert Paul for helpful 
comments, and the NYU Genomics Core facility for sequencing and 
flow cytometry services.

\nolinenumbers

% Either type in your references using
% \begin{thebibliography}{}
% \bibitem{}
% Text
% \end{thebibliography}
%
% or
%
% Compile your BiBTeX database using our plos2015.bst
% style file and paste the contents of your .bbl file
% here. See http://journals.plos.org/plosone/s/latex for 
% step-by-step instructions.
%

\begin{thebibliography}{200}

\bibitem{Hargrove1989}
Hargrove JL, Schmidt FH.
\newblock The role of {mRNA} and protein stability in gene expression.
\newblock FASEB J. 1989;3(12):2360--2370.

\bibitem{Perez-Ortin2013}
P{\'e}rez-Ort{\'\i}n JE, Alepuz P, Ch{\'a}vez S, Choder M.
\newblock Eukaryotic {mRNA} decay: methodologies, pathways, and links to other
  stages of gene expression.
\newblock J Mol Biol. 2013;425(20):3750--3775.

\bibitem{Alonso2012}
Alonso CR.
\newblock A complex {'mRNA} degradation code' controls gene expression during
  animal development.
\newblock Trends Genet. 2012;28(2):78--88.

\bibitem{West2017}
West SM, Mecenas D, Gutwein M, Aristiz{\'a}bal-Corrales D, Piano F, Gunsalus
  KC.
\newblock Developmental dynamics of gene expression and alternative
  polyadenylation in the Caenorhabditis elegans germline.
\newblock Genome Biology. 2018;19(1):8.

\bibitem{Aghib1990}
Aghib D, Bishop J, Ottolenghi S, Guerrasio A, Serra A, Saglio G.
\newblock A 3' truncation of MYC caused by chromosomal translocation in a human
  T-cell leukemia increases mRNA stability.
\newblock Oncogene. 1990;5(5):707—711.

\bibitem{Canadell2015}
Canadell D, Garc{\'\i}a-Mart{\'\i}nez J, Alepuz P, P{\'e}rez-Ort{\'\i}n JE,
  Ari{\~n}o J.
\newblock Impact of high pH stress on yeast gene expression: A comprehensive
  analysis of {mRNA} turnover during stress responses.
\newblock Biochim Biophys Acta. 2015;1849(6):653--664.

\bibitem{Garcia-Martinez2016}
Garc{\'\i}a-Mart{\'\i}nez J, Troul{\'e} K, Ch{\'a}vez S, P{\'e}rez-Ort{\'\i}n
  JE.
\newblock Growth rate controls {mRNA} turnover in steady and non-steady states.
\newblock RNA Biol. 2016;13(12):1175--1181.

\bibitem{Scheffler1998}
Scheffler IE, de~la Cruz BJ, Prieto S.
\newblock Control of {mRNA} turnover as a mechanism of glucose repression in
  Saccharomyces cerevisiae.
\newblock Int J Biochem Cell Biol. 1998;30(11):1175--1193.

\bibitem{Gasch2000}
Gasch AP, Spellman PT, Kao CM, Carmel-Harel O, Eisen MB, Storz G, et~al.
\newblock Genomic expression programs in the response of yeast cells to
  environmental changes.
\newblock Mol Biol Cell. 2000;11(12):4241--4257.

\bibitem{Conway2012}
Conway MK, Grunwald D, Heideman W.
\newblock Glucose, nitrogen, and phosphate repletion in Saccharomyces
  cerevisiae: common transcriptional responses to different nutrient signals.
\newblock G3. 2012;2(9):1003--1017.

\bibitem{Castells2011}
Castells-Roca L, Garc{\'\i}a-Mart{\'\i}nez J, Moreno J, Herrero E, Bell{\'\i}
  G, P{\'e}rez-Ort{\'\i}n JE.
\newblock Heat shock response in yeast involves changes in both transcription
  rates and mRNA stabilities.
\newblock PloS one. 2011;6(2):e17272.

\bibitem{Romero2009}
Romero-Santacreu L, Moreno J, P{\'e}rez-Ort{\'\i}n JE, Alepuz P.
\newblock Specific and global regulation of mRNA stability during osmotic
  stress in Saccharomyces cerevisiae.
\newblock Rna. 2009;15(6):1110--1120.

\bibitem{Molina2008}
Molina-Navarro MM, Castells-Roca L, Bell{\'\i} G, Garc{\'\i}a-Mart{\'\i}nez J,
  Mar{\'\i}n-Navarro J, Moreno J, et~al.
\newblock Comprehensive transcriptional analysis of the oxidative response in
  yeast.
\newblock Journal of Biological Chemistry. 2008;283(26):17908--17918.

\bibitem{Shalem2008}
Shalem O, Dahan O, Levo M, Martinez MR, Furman I, Segal E, et~al.
\newblock Transient transcriptional responses to stress are generated by
  opposing effects of {mRNA} production and degradation.
\newblock Mol Syst Biol. 2008;4:223.

\bibitem{Yin2003}
Yin Z, Wilson S, Hauser NC, Tournu H, Hoheisel JD, Brown AJP.
\newblock Glucose triggers different global responses in yeast, depending on
  the strength of the signal, and transiently stabilizes ribosomal protein
  {mRNAs}.
\newblock Mol Microbiol. 2003;48(3):713--724.

\bibitem{De_la_Cruz2002}
de~la Cruz BJ, Prieto S, Scheffler IE.
\newblock The role of the 5' untranslated region ({UTR}) in glucose-dependent
  {mRNA} decay.
\newblock Yeast. 2002;19(10):887--902.

\bibitem{Mercado1994}
Mercado JJ, Smith R, Sagliocco FA, Brown AJ, Gancedo JM.
\newblock The levels of yeast gluconeogenic {mRNAs} respond to environmental
  factors.
\newblock Eur J Biochem. 1994;224(2):473--481.

\bibitem{Lee2011}
Lee MV, Topper SE, Hubler SL, Hose J, Wenger CD, Coon JJ, et~al.
\newblock A dynamic model of proteome changes reveals new roles for transcript
  alteration in yeast.
\newblock Mol Syst Biol. 2011;7:514.

\bibitem{Kresnowati2006}
Kresnowati MTAP, van Winden WA, Almering MJH, ten Pierick A, Ras C, Knijnenburg
  TA, et~al.
\newblock When transcriptome meets metabolome: fast cellular responses of yeast
  to sudden relief of glucose limitation.
\newblock Mol Syst Biol. 2006;2:49.

\bibitem{Kief1981}
Kief DR, Warner JR.
\newblock Coordinate control of syntheses of ribosomal ribonucleic acid and
  ribosomal proteins during nutritional shift-up in Saccharomyces cerevisiae.
\newblock Mol Cell Biol. 1981;1(11):1007--1015.

\bibitem{Giordano2016}
Giordano N, Mairet F, Gouz{\'e} JL, Geiselmann J, de~Jong H.
\newblock Dynamical Allocation of Cellular Resources as an Optimal Control
  Problem: Novel Insights into Microbial Growth Strategies.
\newblock PLoS Comput Biol. 2016;12(3):e1004802.

\bibitem{Shachrai2010}
Shachrai I, Zaslaver A, Alon U, Dekel E.
\newblock Cost of Unneeded Proteins in E. coli Is Reduced after Several
  Generations in Exponential Growth.
\newblock Mol Cell. 2010;38(5):758--767.

\bibitem{Magasanik2002}
Magasanik B, Kaiser CA.
\newblock Nitrogen regulation in Saccharomyces cerevisiae.
\newblock Gene. 2002;290(1-2):1--18.

\bibitem{Godard2007}
Godard P, Urrestarazu A, Vissers S, Kontos K, Bontempi G, van Helden J, et~al.
\newblock Effect of 21 different nitrogen sources on global gene expression in
  the yeast Saccharomyces cerevisiae.
\newblock Mol Cell Biol. 2007;27(8):3065--3086.

\bibitem{Airoldi2016}
Airoldi EM, Miller D, Athanasiadou R, Brandt N, Abdul-Rahman F, Neymotin B,
  et~al.
\newblock Steady-state and dynamic gene expression programs in Saccharomyces
  cerevisiae in response to variation in environmental nitrogen.
\newblock Mol Biol Cell. 2016;27(8):1383--1396.

\bibitem{Beck1999}
Beck T, Hall MN.
\newblock The {TOR} signalling pathway controls nuclear localization of
  nutrient-regulated transcription factors.
\newblock Nature. 1999;402(6762):689--692.

\bibitem{Tate2013}
Tate JJ, Cooper TG.
\newblock Five conditions commonly used to down-regulate tor complex 1 generate
  different physiological situations exhibiting distinct requirements and
  outcomes.
\newblock J Biol Chem. 2013;288(38):27243--27262.

\bibitem{Tate2017}
Tate JJ, Buford D, Rai R, Cooper TG.
\newblock General Amino Acid Control and 14-3-3 Proteins Bmh1/2 Are Required
  for Nitrogen Catabolite {Repression-Sensitive} Regulation of Gln3 and Gat1
  Localization.
\newblock Genetics. 2017;205(2):633--655.

\bibitem{Cooper1983}
Cooper TG, Sumrada RA.
\newblock What is the function of nitrogen catabolite repression in
  Saccharomyces cerevisiae?
\newblock J Bacteriol. 1983;155(2):623--627.

\bibitem{Stanbrough1995}
Stanbrough M, Magasanik B.
\newblock Transcriptional and posttranslational regulation of the general amino
  acid permease of Saccharomyces cerevisiae.
\newblock J Bacteriol. 1995;177(1):94--102.

\bibitem{Risinger2006}
Risinger AL, Cain NE, Chen EJ, Kaiser CA.
\newblock Activity-dependent reversible inactivation of the general amino acid
  permease.
\newblock Mol Biol Cell. 2006;17(10):4411--4419.

\bibitem{Merhi2012}
Merhi A, Andre B.
\newblock Internal Amino Acids Promote Gap1 Permease Ubiquitylation via
  {TORC1/Npr1/14-3-3-Dependent} Control of the Bul {Arrestin-Like} Adaptors.
\newblock Mol Cell Biol. 2012;32(22):4510--4522.

\bibitem{Parker2012}
Parker R.
\newblock {RNA} Degradation in Saccharomyces cerevisae.
\newblock Genetics. 2012;191(3):671--702.

\bibitem{Tharun2000}
Tharun S, He W, Mayes AE, Lennertz P, Beggs JD, Parker R.
\newblock Yeast Sm-like proteins function in {mRNA} decapping and decay.
\newblock Nature. 2000;404(6777):515--518.

\bibitem{Sharif2013}
Sharif H, Conti E.
\newblock Architecture of the {Lsm1-7-Pat1} complex: a conserved assembly in
  eukaryotic {mRNA} turnover.
\newblock Cell Rep. 2013;5(2):283--291.

\bibitem{Coller2004}
Coller J, Parker R.
\newblock Eukaryotic {mRNA} Decapping.
\newblock Annu Rev Biochem. 2004;73(1):861--890.

\bibitem{Nissan2010}
Nissan T, Rajyaguru P, She M, Song H, Parker R.
\newblock Decapping activators in Saccharomyces cerevisiae act by multiple
  mechanisms.
\newblock Mol Cell. 2010;39(5):773--783.

\bibitem{Olivas2000}
Olivas W, Parker R.
\newblock The Puf3 protein is a transcript-specific regulator of {mRNA}
  degradation in yeast.
\newblock EMBO J. 2000;19(23):6602--6611.

\bibitem{lee2015}
Lee CD, Tu BP.
\newblock Glucose-regulated phosphorylation of the PUF protein Puf3 regulates
  the translational fate of its bound mRNAs and association with RNA granules.
\newblock Cell reports. 2015;11(10):1638--1650.

\bibitem{Haimovich2013}
Haimovich G, Medina DA, Causse SZ, Garber M, Mill{\'a}n-Zambrano G, Barkai O,
  et~al.
\newblock Gene expression is circular: factors for {mRNA} degradation also
  foster {mRNA} synthesis.
\newblock Cell. 2013;153(5):1000--1011.

\bibitem{Trcek2011}
Trcek T, Larson DR, Mold{\'o}n A, Query CC, Singer RH.
\newblock Single-molecule {mRNA} decay measurements reveal promoter- regulated
  {mRNA} stability in yeast.
\newblock Cell. 2011;147(7):1484--1497.

\bibitem{braun2016snf1}
Braun KA, Dombek KM, Young ET.
\newblock Snf1-dependent transcription confers glucose-induced decay upon the
  mRNA product.
\newblock Molecular and cellular biology. 2016;36(4):628--644.

\bibitem{Young2012}
Young ET, Zhang C, Shokat KM, Parua PK, Braun KA.
\newblock The {AMP-activated} protein kinase Snf1 regulates transcription
  factor binding, {RNA} polymerase {II} activity, and {mRNA} stability of
  glucose-repressed genes in Saccharomyces cerevisiae.
\newblock J Biol Chem. 2012;287(34):29021--29034.

\bibitem{Braun2014}
Braun KA, Vaga S, Dombek KM, Fang F, Palmisano S, Aebersold R, et~al.
\newblock Phosphoproteomic analysis identifies proteins involved in
  transcription-coupled {mRNA} decay as targets of Snf1 signaling.
\newblock Sci Signal. 2014;7(333):ra64.

\bibitem{Ramachandran2011}
Ramachandran V, Shah KH, Herman PK.
\newblock The {cAMP-dependent} protein kinase signaling pathway is a key
  regulator of {P} body foci formation.
\newblock Mol Cell. 2011;43(6):973--981.

\bibitem{Luo2011}
Luo G, Costanzo M, Boone C, Dickson RC.
\newblock Nutrients and the Pkh1/2 and Pkc1 protein kinases control {mRNA}
  decay and P-body assembly in yeast.
\newblock J Biol Chem. 2011;286(11):8759--8770.

\bibitem{Talarek2010}
Talarek N, Cameroni E, Jaquenoud M, Luo X, Bontron S, Lippman S, et~al.
\newblock Initiation of the {TORC1-regulated} {G0} program requires Igo1/2,
  which license specific {mRNAs} to evade degradation via the 5'-3' {mRNA}
  decay pathway.
\newblock Mol Cell. 2010;38(3):345--355.

\bibitem{Neymotin2014}
Neymotin B, Athanasiadou R, Gresham D.
\newblock Determination of in vivo {RNA} kinetics using {RATE-seq}.
\newblock RNA. 2014;20(10):1645--1652.

\bibitem{Munchel2011}
Munchel SE, Shultzaberger RK, Takizawa N, Weis K.
\newblock Dynamic profiling of {mRNA} turnover reveals gene-specific and
  system-wide regulation of {mRNA} decay.
\newblock Mol Biol Cell. 2011;22(15):2787--2795.

\bibitem{volland1994}
Volland C, Urban-Grimal D, Geraud G, Haguenauer-Tsapis R.
\newblock Endocytosis and degradation of the yeast uracil permease under
  adverse conditions.
\newblock Journal of Biological Chemistry. 1994;269(13):9833--9841.

\bibitem{Hein1995}
Hein C, Springael JY, Volland C, Haguenauer-Tsapis R, Andre B.
\newblock {NPI1}, an essential yeast gene involved in induced degradation of
  Gap1 and Fur4 permeases, encodes the Rsp5 ubiquitin-protein ligase.
\newblock Mol Microbiol. 1995;18(1):77--87.

\bibitem{Storey2003}
Storey JD, Tibshirani R.
\newblock Statistical significance for genomewide studies.
\newblock Proceedings of the National Academy of Sciences.
  2003;100(16):9440--9445.

\bibitem{Miller2011}
Miller C, Schwalb B, Maier K, Schulz D, D{\"u}mcke S, Zacher B, et~al.
\newblock Dynamic transcriptome analysis measures rates of {mRNA} synthesis and
  decay in yeast.
\newblock Mol Syst Biol. 2011;7:458.

\bibitem{Chen1998}
Chen S, Hyman LE.
\newblock A specific RNA-protein interaction at yeast polyadenylation
  efficiency elements.
\newblock Nucleic acids research. 1998;26(21):4965--4974.

\bibitem{Kessler1997}
Kessler MM, Henry MF, Shen E, Zhao J, Gross S, Silver PA, et~al.
\newblock Hrp1, a sequence-specific RNA-binding protein that shuttles between
  the nucleus and the cytoplasm, is required for mRNA 3'-end formation in
  yeast.
\newblock Genes \& development. 1997;11(19):2545--2556.

\bibitem{Kebaara2003}
Kebaara B.
\newblock The Upf-dependent decay of wild-type {PPR1} {mRNA} depends on its
  {5'-UTR} and first 92 {ORF} nucleotides.
\newblock Nucleic Acids Res. 2003;31(12):3157--3165.

\bibitem{Guisbert2005}
Guisbert K, Duncan K, Li H, Guthrie C.
\newblock Functional specificity of shuttling {hnRNPs} revealed by genome-wide
  analysis of their {RNA} binding profiles.
\newblock RNA. 2005;11(4):383--393.

\bibitem{Khong2017}
Khong A, Matheny T, Jain S, Mitchell SF, Wheeler JR, Parker R.
\newblock The Stress Granule Transcriptome Reveals Principles of {mRNA}
  Accumulation in Stress Granules.
\newblock Mol Cell. 2017;68(4):808--820.e5.

\bibitem{Neklesa2009}
Neklesa TK, Davis RW.
\newblock A genome-wide screen for regulators of {TORC1} in response to amino
  acid starvation reveals a conserved Npr2/3 complex.
\newblock PLoS Genet. 2009;5(6):e1000515.

\bibitem{Sliva2016}
Sliva A, Kuang Z, Meluh PB, Boeke JD.
\newblock Barcode Sequencing Screen Identifies {SUB1} as a Regulator of Yeast
  Pheromone Inducible Genes.
\newblock G3. 2016;6(4):881--892.

\bibitem{Worley2015}
Worley J, Sullivan A, Luo X, Kaplan ME, Capaldi AP.
\newblock {Genome-Wide} Analysis of the {TORC1} and Osmotic Stress Signaling
  Network in Saccharomyces cerevisiae.
\newblock G3. 2015;6(2):463--474.

\bibitem{Vandersluis2014}
VanderSluis B, Hess DC, Pesyna C, Krumholz EW, Syed T, Szappanos B, et~al.
\newblock Broad metabolic sensitivity profiling of a prototrophic yeast
  deletion collection.
\newblock Genome Biol. 2014;15(4):R64.

\bibitem{Smith2009}
Smith AM, Heisler LE, Mellor J, Kaper F, Thompson MJ, Chee M, et~al.
\newblock Quantitative phenotyping via deep barcode sequencing.
\newblock Genome Res. 2009;19(10):1836--1842.

\bibitem{Robinson2013}
Robinson DG, Chen W, Storey JD, Gresham D.
\newblock Design and Analysis of Bar-seq Experiments.
\newblock G3: Genes|Genomes|Genetics. 2013;4(1):11--18.

\bibitem{Giaever2014}
Giaever G, Nislow C.
\newblock The yeast deletion collection: a decade of functional genomics.
\newblock Genetics. 2014;197(2):451--465.

\bibitem{Kinney2010}
Kinney JB, Murugan A, Callan CG, Cox EC.
\newblock Using deep sequencing to characterize the biophysical mechanism of a
  transcriptional regulatory sequence.
\newblock Proceedings of the National Academy of Sciences.
  2010;107(20):9158--9163.

\bibitem{Peterman2016}
Peterman N, Levine E.
\newblock Sort-seq under the hood: implications of design choices on
  large-scale characterization of sequence-function relations.
\newblock BMC Genomics. 2016;17:206.

\bibitem{Raj2008}
Raj A, van~den Bogaard P, Rifkin SA, van Oudenaarden A, Tyagi S.
\newblock Imaging individual {mRNA} molecules using multiple singly labeled
  probes.
\newblock Nat Methods. 2008;5(10):877--879.

\bibitem{Klemm2014}
Klemm S, Semrau S, Wiebrands K, Mooijman D, Faddah DA, Jaenisch R, et~al.
\newblock Transcriptional profiling of cells sorted by {RNA} abundance.
\newblock Nat Methods. 2014;11(5):549--551.

\bibitem{Hanley2013}
Hanley MB, Lomas W, Mittar D, Maino V, Park E.
\newblock Detection of low abundance {RNA} molecules in individual cells by
  flow cytometry.
\newblock PLoS One. 2013;8(2):e57002.

\bibitem{Fu2011}
Fu GK, Hu J, Wang PH, Fodor SPA.
\newblock Counting individual {DNA} molecules by the stochastic attachment of
  diverse labels.
\newblock Proc Natl Acad Sci U S A. 2011;108(22):9026--9031.

\bibitem{Scherens2006}
Scherens B, Feller A, Vierendeels F, Messenguy F, Dubois E.
\newblock Identification of direct and indirect targets of the Gln3 and Gat1
  activators by transcriptional profiling in response to nitrogen availability
  in the short and long term.
\newblock FEMS Yeast Res. 2006;6(5):777--791.

\bibitem{Hoyle2007}
Hoyle NP, Castelli LM, Campbell SG, Holmes LEA, Ashe MP.
\newblock Stress-dependent relocalization of translationally primed {mRNPs} to
  cytoplasmic granules that are kinetically and spatially distinct from
  P-bodies.
\newblock J Cell Biol. 2007;179(1):65--74.

\bibitem{Rajyaguru2012}
Rajyaguru P, She M, Parker R.
\newblock Scd6 targets {eIF4G} to repress translation: {RGG} motif proteins as
  a class of {eIF4G-binding} proteins.
\newblock Mol Cell. 2012;45(2):244--254.

\bibitem{Puig2005}
Puig S, Askeland E, Thiele DJ.
\newblock Coordinated remodeling of cellular metabolism during iron deficiency
  through targeted {mRNA} degradation.
\newblock Cell. 2005;120(1):99--110.

\bibitem{Duffy2015}
Duffy EE, Rutenberg-Schoenberg M, Stark CD, Kitchen RR, Gerstein MB, Simon MD.
\newblock Tracking Distinct {RNA} Populations Using Efficient and Reversible
  Covalent Chemistry.
\newblock Mol Cell. 2015;59(5):858--866.

\bibitem{Chan2017}
Chan LY, Mugler CF, Heinrich S, Vallotton P, Weis K.
\newblock Non-invasive measurement of mRNA decay reveals translation initiation
  as the major determinant of mRNA stability.
\newblock bioRxiv. 2017; p. 214775.

\bibitem{Boles1998}
Boles E, de~Jong-Gubbels P, Pronk JT.
\newblock Identification and Characterization of MAE1, the Saccharomyces
  cerevisiae Structural Gene Encoding Mitochondrial Malic Enzyme.
\newblock Journal of bacteriology. 1998;180(11):2875--2882.

\bibitem{Tesniere2017}
Tesni{\`e}re C, Pradal M, Bessi{\`e}re C, Sanchez I, Blondin B, Bigey F.
\newblock Relief from nitrogen starvation triggers a transient destabilization
  of glycolytic {mRNAs} in Saccharomyces cerevisiae cells.
\newblock Mol Biol Cell. 2017;.

\bibitem{Gonzalez2000}
Gonz{\'a}lez CI, Ruiz-Echevarr{\'\i}a MJ, Vasudevan S, Henry MF, Peltz SW.
\newblock The yeast {hnRNP-like} protein {Hrp1/Nab4} marks a transcript for
  nonsense-mediated {mRNA} decay.
\newblock Mol Cell. 2000;5(3):489--499.

\bibitem{Stanbrough1996}
Stanbrough M, Magasanik B.
\newblock Two transcription factors, Gln3p and Nil1p, use the same GATAAG sites
  to activate the expression of GAP1 of Saccharomyces cerevisiae.
\newblock Journal of bacteriology. 1996;178(8):2465--2468.

\bibitem{Arribere2013}
Arribere JA, Gilbert WV.
\newblock Roles for transcript leaders in translation and mRNA decay revealed
  by transcript leader sequencing.
\newblock Genome research. 2013;23(6):977--987.

\bibitem{Baumgartner2011}
Baumgartner BL, Bennett MR, Ferry M, Johnson TL, Tsimring LS, Hasty J.
\newblock Antagonistic gene transcripts regulate adaptation to new growth
  environments.
\newblock Proc Natl Acad Sci U S A. 2011;108(52):21087--21092.

\bibitem{Neymotin2016}
Neymotin B, Ettore V, Gresham D.
\newblock Multiple Transcript Properties Related to Translation Affect {mRNA}
  Degradation Rates in Saccharomyces cerevisiae.
\newblock G3. 2016;.

\bibitem{Rouhanifard2017}
Rouhanifard SH, Dunagin M, Mellis IA, Bayatpour S, Symmons O, Cote A, et~al.
\newblock Single-molecule fluorescent amplification of {RNA} using {clampFISH}
  probes.
\newblock bioRxiv (pre-print). 2017;.

\bibitem{Brauer2008}
Brauer MJ, Huttenhower C, Airoldi EM, Rosenstein R, Matese JC, Gresham D,
  et~al.
\newblock Coordination of growth rate, cell cycle, stress response, and
  metabolic activity in yeast.
\newblock Molecular biology of the cell. 2008;19(1):352--367.

\bibitem{Amberg2005}
Amberg DC, Burke DJ, Strathern JN.
\newblock Methods in Yeast Genetics: A Cold Spring Harbor Laboratory Course
  Manual, 2005 Edition (Cold Spring).
\newblock $\{$Cold Spring Harbor Laboratory Press$\}$; 2005.

\bibitem{Storici2006}
Storici F, Resnick MA.
\newblock The delitto perfetto approach to in vivo site-directed mutagenesis
  and chromosome rearrangements with synthetic oligonucleotides in yeast.
\newblock Methods in enzymology. 2006;409:329--345.

\bibitem{Stacklies2007}
Stacklies W, Redestig H, Scholz M, Walther D, Selbig J.
\newblock pcaMethods a bioconductor package providing {PCA} methods for
  incomplete data.
\newblock Bioinformatics. 2007;23(9):1164--1167.

\bibitem{Hong2017}
Hong J, Gresham D.
\newblock Incorporation of unique molecular identifiers in TruSeq adapters
  improves the accuracy of quantitative sequencing.
\newblock BioTechniques. 2017;63(5):221--226.

\bibitem{Martin2011}
Martin M.
\newblock Cutadapt removes adapter sequences from high-throughput sequencing
  reads.
\newblock EMBnet journal. 2011;17(1):pp--10.

\bibitem{Kim2013}
Kim D, Pertea G, Trapnell C, Pimentel H, Kelley R, Salzberg SL.
\newblock {TopHat2}: accurate alignment of transcriptomes in the presence of
  insertions, deletions and gene fusions.
\newblock Genome Biol. 2013;14(4):R36.

\bibitem{Li2009}
Li H, Handsaker B, Wysoker A, Fennell T, Ruan J, Homer N, et~al.
\newblock The sequence alignment/map format and SAMtools.
\newblock Bioinformatics. 2009;25(16):2078--2079.

\bibitem{Smith2017}
Smith T, Heger A, Sudbery I.
\newblock {UMI-tools}: modeling sequencing errors in Unique Molecular
  Identifiers to improve quantification accuracy.
\newblock Genome Res. 2017;27(3):491--499.

\bibitem{Anders2015}
Anders S, Pyl PT, Huber W.
\newblock HTSeq—a Python framework to work with high-throughput sequencing
  data.
\newblock Bioinformatics. 2015;31(2):166--169.

\bibitem{Storey2015}
Storey JD, Bass AJ, Dabney A, Robinson D. qvalue: Q-value estimation for false
  discovery rate control; 2015.
\newblock Available from: \url{http://github.com/jdstorey/qvalue}.

\bibitem{Huggins2011}
Huggins P, Zhong S, Shiff I, Beckerman R, Laptenko O, Prives C, et~al.
\newblock {DECOD}: fast and accurate discriminative {DNA} motif finding.
\newblock Bioinformatics. 2011;27(17):2361--2367.

\bibitem{Elemento2007}
Elemento O, Slonim N, Tavazoie S.
\newblock A universal framework for regulatory element discovery across all
  genomes and data types.
\newblock Mol Cell. 2007;28(2):337--350.

\bibitem{Goodarzi2012}
Goodarzi H, Najafabadi HS, Oikonomou P, Greco TM, Fish L, Salavati R, et~al.
\newblock Systematic discovery of structural elements governing stability of
  mammalian messenger {RNAs}.
\newblock Nature. 2012;485(7397):264--268.

\bibitem{Li2010}
Li X, Quon G, Lipshitz HD, Morris Q.
\newblock Predicting in vivo binding sites of {RNA-binding} proteins using
  {mRNA} secondary structure.
\newblock RNA. 2010;16(6):1096--1107.

\bibitem{Ray2013}
Ray D, Kazan H, Cook KB, Weirauch MT, Najafabadi HS, Li X, et~al.
\newblock A compendium of {RNA-binding} motifs for decoding gene regulation.
\newblock Nature. 2013;499(7457):172--177.

\bibitem{McLeay2010}
McLeay RC, Bailey TL.
\newblock Motif Enrichment Analysis: a unified framework and an evaluation on
  {ChIP} data.
\newblock BMC Bioinformatics. 2010;11:165.

\bibitem{Lawrence2013}
Lawrence M, Huber W, Pages H, Aboyoun P, Carlson M, Gentleman R, et~al.
\newblock Software for computing and annotating genomic ranges.
\newblock PLoS computational biology. 2013;9(8):e1003118.

\bibitem{Pelechano2014}
Pelechano V, Wei W, Jakob P, Steinmetz LM.
\newblock Genome-wide identification of transcript start and end sites by
  transcript isoform sequencing.
\newblock Nat Protoc. 2014;9(7):1740--1759.

\bibitem{Freeberg2013}
Freeberg MA, Han T, Moresco JJ, Kong A, Yang YC, Lu ZJ, et~al.
\newblock Pervasive and dynamic protein binding sites of the {mRNA}
  transcriptome in Saccharomyces cerevisiae.
\newblock Genome Biol. 2013;14(2):R13.

\bibitem{Mcisaac2013}
McIsaac RS, Silverman SJ, Parsons L, Xu P, Briehof R, McClean MN, et~al.
\newblock Visualization and analysis of mRNA molecules using fluorescence in
  situ hybridization in Saccharomyces cerevisiae.
\newblock Journal of visualized experiments: JoVE. 2013;(76).

\bibitem{Cock2009}
Cock PJA, Antao T, Chang JT, Chapman BA, Cox CJ, Dalke A, et~al.
\newblock Biopython: freely available Python tools for computational molecular
  biology and bioinformatics.
\newblock Bioinformatics. 2009;25(11):1422--1423.

\bibitem{Li2013}
Li H.
\newblock Aligning sequence reads, clone sequences and assembly contigs with
  {BWA-MEM}.
\newblock arXiv. 2013; p. 1303.3997.

\bibitem{Team2000}
Team RC.
\newblock R language definition.
\newblock Vienna, Austria: R foundation for statistical computing. 2000;.

\bibitem{Yu2012}
Yu G, Wang LG, Han Y, He QY.
\newblock clusterProfiler: an R package for comparing biological themes among
  gene clusters.
\newblock Omics: a journal of integrative biology. 2012;16(5):284--287.

\bibitem{Presnyak2015}
Presnyak V, Alhusaini N, Chen YH, Martin S, Morris N, Kline N, et~al.
\newblock Codon optimality is a major determinant of {mRNA} stability.
\newblock Cell. 2015;160(6):1111--1124.

\end{thebibliography}

%\bibliographystyle{scripts/plos2015}
%\bibliography{scripts/miller2018_references}

% Place figure captions after the first paragraph in which they are cited.
%\begin{figure}[!h]
%\caption{{\bf Bold the figure title.}
%Figure caption text here, please use this space for the figure panel descriptions instead of using subfigure commands. A: Lorem ipsum dolor sit amet. B: Consectetur adipiscing elit.}
%\label{fig1}
%\end{figure}
% Place tables after the first paragraph in which they are cited.
%\begin{table}[!ht]
%\begin{adjustwidth}{-2.25in}{0in} % Comment out/remove adjustwidth environment if table fits in text column.
%\centering
%\caption{
%{\bf Table caption Nulla mi mi, venenatis sed ipsum varius, volutpat euismod diam.}}
%\begin{tabular}{|l+l|l|l|l|l|l|l|}
%\hline
%\multicolumn{4}{|l|}{\bf Heading1} & \multicolumn{4}{|l|}{\bf Heading2}\\ \thickhline
%$cell1 row1$ & cell2 row 1 & cell3 row 1 & cell4 row 1 & cell5 row 1 & cell6 row 1 & cell7 row 1 & cell8 row 1\\ \hline
%$cell1 row2$ & cell2 row 2 & cell3 row 2 & cell4 row 2 & cell5 row 2 & cell6 row 2 & cell7 row 2 & cell8 row 2\\ \hline
%$cell1 row3$ & cell2 row 3 & cell3 row 3 & cell4 row 3 & cell5 row 3 & cell6 row 3 & cell7 row 3 & cell8 row 3\\ \hline
%\end{tabular}
%\begin{flushleft} Table notes Phasellus venenatis, tortor nec vestibulum mattis, massa tortor interdum felis, nec pellentesque metus tortor nec nisl. Ut ornare mauris tellus, vel dapibus arcu suscipit sed.
%\end{flushleft}
%\label{table1}
%\end{adjustwidth}
%\end{table}

\end{document}

